\documentclass[11pt]{article}
\usepackage[T1]{fontenc}
\usepackage{lmodern}
\usepackage[margin=1in]{geometry}
\usepackage{amsmath,amssymb,amsthm,mathtools}
\usepackage{microtype}
\usepackage{booktabs,longtable,array}
\usepackage[dvipsnames]{xcolor}
\usepackage[hyphens]{url}
\usepackage{hyperref}
\hypersetup{colorlinks=true,linkcolor=MidnightBlue,citecolor=MidnightBlue,
 urlcolor=MidnightBlue,pdftitle={Bell Numbers in Intuitionistic Primitive Recursive Arithmetic},
 pdfauthor={Research prepared for Vladimir Reshetnikov},
 pdfsubject={Constructive Bell number proofs, effective asymptotics, and Lean formalization blueprint}}
\newcommand{\N}{\mathbb N}
\newcommand{\Z}{\mathbb Z}
\newcommand{\Q}{\mathbb Q}
\newcommand{\R}{\mathbb R}
\newcommand{\ipra}{\ensuremath{\mathrm{iPRA}}}
\pdfstringdefDisableCommands{\def\ipra{iPRA}}
\newcommand{\Bell}{B}
\newtheorem{theorem}{Theorem}[section]
\newtheorem{lemma}[theorem]{Lemma}
\newtheorem{proposition}[theorem]{Proposition}
\newtheorem{corollary}[theorem]{Corollary}
\theoremstyle{definition}
\newtheorem{definition}[theorem]{Definition}
\theoremstyle{remark}
\newtheorem{remark}[theorem]{Remark}
\numberwithin{equation}{section}
\setlength{\emergencystretch}{2em}
\setlength{\parskip}{3pt}
\setcounter{tocdepth}{2}
\allowdisplaybreaks[2]

\begin{document}
\hypersetup{pageanchor=false}
\begin{titlepage}
\centering
\vspace*{1.2cm}
{\Large Constructive arithmetic research\par}
\vspace{1.1cm}
{\huge\bfseries Bell Numbers in Intuitionistic\\[8pt]
Primitive Recursive Arithmetic\par}
\vspace{0.8cm}
{\Large Finite proofs, effective asymptotics,\\[5pt]
and a Lean formalization blueprint\par}
\vspace{1.1cm}
{\large Prepared for Vladimir Reshetnikov\par}
\vspace{0.3cm}
{\large November 15, 2026\par}
\vfill
\begin{minipage}{0.88\textwidth}
\textbf{Principal result.} The standard Bell identities and the principal
growth formulas on the two requested pages admit proofs with finite
primitive recursive content. We supply the combinatorial proofs, effective
analytic estimates, both logarithmic-convexity inequalities, modular periods,
the universal $0.792$ upper bound, and the two-correction saddle expansion.

\medskip
\textbf{Formal status.} A constructive Lean core is compiled and audited.
The remaining targets have explicit arithmetic definitions and a detailed
route through the inspected ProveIt proof calculus. The analytic theorems
are proved here in human-readable arithmetic form; the package does not
claim that their full Lean or object-language derivations have already
been completed.
\end{minipage}
\vfill
{\small ProveIt source snapshot:\quad
\texttt{c0d0dc37f60aea10c0182ba4eece9eae8c7e13c6}\par}
\end{titlepage}
\hypersetup{pageanchor=true}
\pagenumbering{roman}

\begin{abstract}
We develop the Bell numbers in intuitionistic primitive recursive
arithmetic, with infinite expressions interpreted by specified rational
approximants and primitive recursive error bounds. The finite development
includes partition and permutation interpretations, all principal summation
identities on the requested pages, a Charlier-polynomial Hankel determinant
proof, strict log-convexity, normalized log-concavity, Touchard congruences,
and constructive periods modulo every integer. The analytic development
derives Dobi\'nski's formula, its ceiling truncation, generating functions,
hypergeometric and integral representations, and effective growth results.
In particular, it proves the de Bruijn logarithmic expansion with a numerical
remainder, the Lambert $W$ equivalents with the explicit threshold
$n\ge2^{64q^2}$ for relative error $1/q$, and a two-correction expansion
uniform under logarithmic shifts. The all-$n$ bound with constant $0.792$
uses an elementary estimate for $n\ge512$ and a reproducible exact rational
certificate for the smaller cases. Checked Lean code establishes a finite
core; executable arithmetic specifications and a module-by-module plan
describe the remaining formalization.
\end{abstract}

\tableofcontents
\newpage
\pagenumbering{arabic}

\section*{Reading guide and coverage}
\addcontentsline{toc}{section}{Reading guide and coverage}

The mathematical proofs are in Sections~\ref{sec:foundations}--\ref{sec:corrections}.
Section~\ref{sec:formalization} states exactly what is already Lean-checked
and defines the remaining arithmetic targets. The final numerical appendix
records reproducible ground calculations. Throughout, a formal series means
finite coefficient identities, and an analytic expression means its specified
finite-precision arithmetic interpretation.

\small
\begin{longtable}{@{}>{\raggedright\arraybackslash}p{0.37\textwidth}>{\raggedright\arraybackslash}p{0.49\textwidth}p{0.08\textwidth}@{}}
\toprule
\textbf{Requested family} & \textbf{Result supplied} & \textbf{Section}\\
\midrule
\endfirsthead
\toprule
\textbf{Requested family} & \textbf{Result supplied} & \textbf{Section}\\
\midrule
\endhead
Partitions, equivalence relations, rhyme schemes, squarefree factorizations &
Explicit bounded encodings and inverse maps & \ref{sec:finite}\\
Permutation and shuffle interpretations &
All eight listed generalized avoidance classes, the $321$ extension class,
and return shuffles & \ref{sec:finite}\\
Stirling and binomial recurrences; triangle &
Finite proofs, initial values and Bell triangle & \ref{sec:finite}\\
Finite summation formulas &
Spivey, inclusion--exclusion, derangement form, Pascal and signed-Stirling
identities & \ref{sec:finite}\\
Hankel determinant &
Exact product of factorials via finite orthogonality & \ref{sec:finite}\\
Log-convexity and normalized log-concavity &
Both inequalities, strictness of the first, and product bounds & \ref{sec:finite}\\
Dobi\'nski and Comtet &
Finite truncation identity, explicit tail, and ceiling formula & \ref{sec:series}\\
Poisson moments and Bell polynomials &
Specified weighted sums and finite polynomial identities & \ref{sec:series}\\
Exponential generating function &
Formal and entire identities, including tail estimates & \ref{sec:series}\\
Ordinary generating series &
Zero radius; negative-axis function with full asymptotic remainder & \ref{sec:series}\\
Hypergeometric formulas &
Factorial-series proofs of the Bell and confluent expressions, plus the
alternating rising-factorial series & \ref{sec:series}\\
Integral representations &
Circle coefficient extraction and the corrected Ces\`aro integral & \ref{sec:series}\\
Touchard and prime-power shifts &
Explicit orbit and finite-polynomial proofs & \ref{sec:congruences}\\
Periods &
Prime period, arbitrary-modulus bounded search, and the listed least periods
for moduli $1$ through $12$ & \ref{sec:congruences}\\
de Bruijn expansion &
The printed expansion with an explicit error and threshold & \ref{sec:growth}\\
Lambert $W$ equivalents &
All leading equivalents on the pages, with relative-error moduli & \ref{sec:growth}\\
Moser--Wyman shifted expansion &
All $P_i,Q_i$ displayed and an explicit uniform cubic remainder & \ref{sec:corrections}\\
$0.792$ all-$n$ upper bound &
Proof for $n\ge512$ and exact certificates for $1\le n\le511$ & \ref{sec:universal-bound}\\
Eventual $e^{-0.6+\eta}$ upper bound &
A different explicit threshold, $n\ge2^{16}$, already sufficient with
$\eta=0$ & \ref{sec:universal-bound}\\
Selected numerical data &
Digit counts through $B_{1000}$ and primality checks through index $13$ &
\ref{sec:numeric}\\
\bottomrule
\end{longtable}
\normalsize

\paragraph{Boundaries of the coverage.}
We do not prove the conjecture that infinitely many Bell numbers are prime,
or that the displayed prime period is always least. We do not reproduce the
large primality computations, exhaustive prime searches, or the digit count
for $B_{10000}$. The branch-dependent incomplete-gamma expression is omitted.
The contour formula is established on circles around zero; a general contour
deformation theorem is unnecessary for the Bell growth proof and is not
claimed. The exact inverse-function threshold printed for the eventual
$e^{-0.6+\eta}$ bound is replaced by the explicit threshold above. The
Stanley--Wilf theorem, a separate general result about ordinary permutation
avoidance, is not re-proved. These exclusions are recorded so that the broad
coverage is not mistaken for a proof of every historical or computational
sentence on the pages.

\section{Scope and the constructive arithmetic setting}\label{sec:foundations}

The Bell numbers provide a particularly tractable test of how much apparent
analysis can be expressed in intuitionistic primitive recursive arithmetic.
Their defining objects are finite, and the analytic formulas used below have
effective, explicitly bounded remainders.  Our aim is to give proofs of the
statements, not merely to observe that their numerical instances can be tested.
The target inventory is the two pages specified in the question
\cite{mathworld,wikipedia}; attribution of an identity to these pages does not
make either page an axiom of our argument.

\subsection{What ``in \ipra'' means here}

We use first-order intuitionistic logic with equality, symbols for primitive
recursive functions, their defining equations, the usual elementary successor
axioms, and induction for quantifier-free formulas.  An equivalent presentation
may use Boolean characteristic functions and induction for equations.  There is
no induction for arbitrary quantified formulas, no excluded middle for
arbitrary propositions, and no quantification over arbitrary functions or real
numbers.  Free numerical parameters are allowed in each induction instance.

The paper has two kinds of assertions.
\begin{enumerate}
\item \emph{Exact finite assertions}: equations, inequalities, divisibility,
finite counting statements and bounded algorithms, written directly using
primitive recursive operations.
\item \emph{Represented analytic assertions}: abbreviations for families of
finite rational inequalities with specified errors and primitive recursive
precision bounds.  For example, equality of two explicitly constructed real
names means that their rational approximants differ by at most the sum of the
declared errors, at every precision.  An asymptotic equivalence includes a
primitive recursive threshold for each rational accuracy request.
\end{enumerate}
This convention does not assert that the first-order theory contains a complete
ordered field of real numbers.  It also does not identify a theorem proved in
Lean's ambient type theory with a syntactic derivation in \ipra.  The distinction
is maintained in Section~\ref{sec:formalization}.

\subsection{Finite objects and bounded induction}

We encode a length-$l$ list of entries smaller than $b$ by its base-$b$ numeral,
with its length recorded separately.  Take $b\ge2$; exceptional smaller bases
are replaced by $2$.  Reading an entry uses division and remainder.  A table
with variable row lengths can instead use nested pairing; either convention is
primitive recursive.  A bounded sum, bounded product, bounded maximum and
bounded search are defined by primitive recursion.  Integers are signed
natural-number codes, and rationals are pairs $(z,d)$ with $d>0$.  Reduction to
lowest terms is optional.  Comparisons are decided by cross multiplication.
Polynomial and matrix computations use finite coefficient and entry lists.

A bounded formula is replaced by its Boolean characteristic function.  For
example, ``all $j\le m$ satisfy $P(j)$'' becomes the equation
\[
 \prod_{j=0}^{m}\chi_P(j)=1.
\]
Thus an induction on a bounded assertion is an instance of quantifier-free
induction after this explicit translation.  Nested induction on fixed finite
tuples is allowed: it is a finite sequence of induction arguments, not induction
along a new transfinite ordering.  Induction on a table of length $m$ is made
simultaneous by coding the whole table invariant in one characteristic function.

For a finite injection or bijection we supply algorithms for the map and its
inverse, together with bounds on their searches.  The counting lemma itself is
proved by induction while scanning the bounded domain.  Likewise, finite
pigeonhole is proved by keeping the Boolean table of already seen values.
These observations justify the finite counting arguments below without using
an infinitary set-existence principle.

\begin{remark}[A limitation of the compilation claim]
It is not true that every computable construction, every terminating Lean
program, or every classical proof of a numerical identity automatically belongs
to \ipra.  In this paper the loops have displayed primitive recursive bounds,
and the invariants used in the arithmetic arguments are bounded.  The prose
proof is an arithmetic proof blueprint; a complete generated object-language
derivation is a separate formalization task, not a deliverable claimed to have
already been checked.
\end{remark}

\subsection{Elementary analytic names}\label{subsec:names}

Only fixed, explicitly given families of names occur.  There is no claimed
primitive recursive evaluator for arbitrary primitive recursive programs.
An error parameter $s\in\N$ means absolute error at most $2^{-s}$.
When a named quantity is supplied as an argument, its data also include a
certified rational range and, for a logarithm or reciprocal, a strictly
positive rational lower bound. One first evaluates the argument to a
predetermined precision and then the rational approximation program. A
derivative bound $M$ on the certified range makes the propagated error at
most $M$ times the input error, plus the program's own error. This explains
composition of the names without exact sign tests on arbitrary real names.

For $|x|\le A$, with $A$ a positive integer, let
\[
 E_M(x)=\sum_{j=0}^{M}\frac{x^j}{j!}.
\]
If $M+2\ge2A$, the absolute tail is bounded by
\begin{equation}\label{eq:exp-tail}
 \sum_{j>M}\frac{A^j}{j!}\le
 \frac{2A^{M+1}}{(M+1)!}.
\end{equation}
The inequality means that every finite tail satisfies the bound; it follows
from a geometric majorant because the successive ratios are at most $1/2$.
One completely explicit precision rule is obtained by starting at $K=2A$:
$A^K/K!\le A^{2A}\le2^{2A^2}$, so
$M=2A+2A^2+s+3$ is more than sufficient.  This also proves convergence and the
addition and multiplication formulas by comparing finite products and their
tails.  Negative exponents may alternatively use the reciprocal of a positive
exponential, whose lower bound is $1$.

For $x>0$ rational, first find, by bounded integer comparisons, $k\in\mathbb Z$
and $y\in[1,2]$ such that $x=2^k y$.  The magnitude of $k$ is bounded by the sum
of the numerator and denominator of $x$.  With $z=(y-1)/(y+1)$, define
\begin{equation}\label{eq:log-series}
 \log y=2\sum_{j=0}^{J-1}\frac{z^{2j+1}}{2j+1}+R_J,
 \qquad
 0\le R_J\le\frac{2z^{2J+1}}{(2J+1)(1-z^2)}.
\end{equation}
Here $0\le z\le1/3$.  The identity is obtained by integrating the finite
geometric identity for $2/(1-z^2)$, retaining its remainder; equivalently it
can be proved by finite differentiation of the displayed polynomials.  Set
$\log x=k\log2+\log y$, using the same formula for $\log2$.  Taking
$J=s+|k|+4$ suffices for the present absolute-error convention.  The derivative
$1/x$, monotonicity, the product law and the inverse relation with $\exp$ follow
by finite polynomial approximation on an interval with an explicit positive
lower endpoint.  The error in replacing such an identity by a rational one is
bounded by the displayed tails and a bound for the derivative on that interval.

Square roots of nonnegative rational numbers are obtained by bisection using
the rational test $q^2\le x$.  Positive $d$th roots use $q^d\le x$ in exactly the
same way.  An initial interval $[0,1+\lceil x\rceil]$ and
$s+\lceil x\rceil+2$ bisections suffice.  Powers $x^\alpha$ at positive named
$x$ mean $\exp(\alpha\log x)$ unless a root construction is simpler.

The sine and cosine names are their factorial series.  The same absolute-tail
estimate as \eqref{eq:exp-tail} applies on a bounded interval.  One may define
$\pi$ by the convergent arctangent construction
$\pi=16\arctan(1/5)-4\arctan(1/239)$, where
\[
 \arctan t=\sum_{j=0}^{J-1}\frac{(-1)^j t^{2j+1}}{2j+1}
             +\rho_J,\qquad |\rho_J|\le\frac{|t|^{2J+1}}{2J+1}.
\]
The addition formula for tangent verifies this normalization; it agrees with
$4\arctan1$, whose alternating partial sums also give a slower primitive
recursive name.  Polynomial approximation proves the addition formulas for
sine and cosine, their derivatives, and $\exp(it)=\cos t+i\sin t$.
In particular, orthogonality of integer Fourier modes follows by integrating
their explicit sine and cosine antiderivatives.

\subsection{Integration and Taylor estimates as finite statements}

For the functions used here, an integral over a rational compact interval is a
name given by equal-mesh Riemann sums.  If $|f'|\le M$ on an interval of length
$L$, the error of a left-endpoint sum with $q$ subintervals is at most
$ML^2/q$.  Indeed, on each subinterval the oscillation is at most its length
times $M$; summing the $q$ resulting rectangles gives the bound.  The usual
derivative bounds used below are obtained directly from factorial series.
Endpoints such as $\pi$ are approximated with a certified error, adding at most
$\sup|f|$ times the endpoint error.  Thus all integrations in the proofs have
primitive recursive mesh bounds.

Taylor's estimate
\begin{equation}\label{eq:taylor-estimate}
 \left|f(x+h)-\sum_{j=0}^{d}\frac{f^{(j)}(x)}{j!}h^j\right|
 \le\frac{M|h|^{d+1}}{(d+1)!}
\end{equation}
holds when $|f^{(d+1)}|\le M$ on the segment.  A constructive proof integrates
the derivative identity repeatedly and bounds the resulting simplex integral;
the simplex has volume $|h|^{d+1}/(d+1)!$, verified by induction.  Equivalently,
one first proves the formula for polynomials and then uses uniform factorial
tails.  No intermediate-value theorem or selection of an unknown Taylor point
is used.  The elementary estimates $1-u\le e^{-u}$,
$|e^z-1|\le e^{|z|}-1$, and $1-\cos t\ge t^2/3$ for $|t|\le1$ follow from
these same finite series estimates.

Unbounded Gaussian integrals are notation for truncated integrals with the
tail certificate
\[
 \int_T^\infty e^{-t^2/2}\,dt\le T^{-1}e^{-T^2/2}\quad(T>0).
\]
Integration by parts gives
\[
 \frac1{\sqrt{2\pi}}\int_{-\infty}^{\infty}t^{2j}e^{-t^2/2}\,dt
 =(2j-1)!!,
 \qquad \int_{-\infty}^{\infty}t^{2j+1}e^{-t^2/2}\,dt=0.
\]
For completeness, the normalizing integral is computed by squaring a truncated
Gaussian integral, sandwiching the square between inscribed and circumscribed
disks, and using finite radial annuli. Integration of
$r e^{-r^2/2}$ gives the limit $2\pi$ for the square, with errors bounded by
Gaussian tails. More explicitly, an inscribed regular $q$-gon has area
$(q/2)R^2\sin(2\pi/q)$; for $q\ge8$, its difference from $\pi R^2$ is at
most $4\pi^3R^2/(3q^2)$ by the cubic sine remainder. The circumscribed
polygon has area $qR^2\tan(\pi/q)$, with excess at most
$2\pi^3R^2/(3q^2)$: use
$\sin x-x\cos x=\int_0^x t\sin t\,dt\le x^3/3$ and
$\cos(\pi/q)\ge1/2$. These two
polygonal computations first establish the disk-area name. On radial steps
$R/M$, the function $e^{-r^2/2}$ oscillates by at most $R^2/M$; the total
annular integration error is at most $4R^4/M$. Taking $q,M$ to be explicit
powers of $2$, with the requested error and $R$ as parameters, therefore
justifies the disk integral calculation by finite polygonal and rectangular
sums. This is an effective calculation, not an appeal to general measure
theory.

Every passage to a limit below can consequently be read as follows: fix a
rational requested error, choose the indicated finite truncations and meshes,
prove the finite rational inequality, and add its explicit error allowances.
The exact integer conclusions are unaffected by this notation.

\section{Finite arithmetic and combinatorics}
\label{sec:finite}

\subsection{A primitive recursive definition}

Write $[n]=\{0,\ldots,n-1\}$. A \emph{restricted growth word} of length $n$ is a word
$a_0\cdots a_{n-1}$ with $a_0=0$ when $n>0$ and
\[
  a_i\leq 1+\max_{j<i}a_j\qquad(0<i<n).
\]
The empty word is allowed. Its number of blocks is zero; that of a nonempty
word is $1+\max_i a_i$. Encode a word by its $n$ base-$(n+1)$ digits, retaining
the length separately. Validity and block count are bounded-loop tests, and
there are at most $(n+1)^n$ candidate codes. Let $S(n,k)$ count valid words
with $k$ blocks, and put
\begin{equation}\label{eq:bell-stirling}
  B_n=\sum_{k=0}^{n}S(n,k).
\end{equation}
In particular $B_0=1$. This definition is primitive recursive without any
unbounded search.

The blocks of a word are the nonempty sets $\{i:a_i=j\}$. Conversely, order the
blocks of a partition by their least elements and give them labels
$0,1,\ldots$. These two algorithms are inverse. Consequently $B_n$ counts
partitions of $[n]$, equivalence relations on $[n]$, and rhyme schemes after
canonical relabelling. An equivalence relation is encoded by an $n$ by $n$
Boolean matrix; its classes are found by scanning for their least members.
This also proves the equivalence-relation interpretation without forming a
quotient of an infinite set.

\begin{proposition}[Stirling recurrence and positivity]\label{prop:stirling}
With $S(n,k)=0$ outside $0\leq k\leq n$,
\[
 S(0,0)=1,\qquad S(n+1,k)=kS(n,k)+S(n,k-1).
\]
Thus $B_n\geq1$, and, for $n\geq1$,
\[
 2^{n-1}\leq B_n\leq n^n.
\]
The first eleven values are
\[
 1,1,2,5,15,52,203,877,4140,21147,115975.
\]
\end{proposition}
\begin{proof}
The last element of a partition either joins one of its $k$ pre-existing
blocks or is the sole member of a new block. Deleting that element gives the
displayed recurrence and its inverse construction. The all-singleton
partition proves positivity. For the lower bound, choose any subset of the
other $n-1$ elements to join the block containing $0$, and make every remaining
element a singleton. These $2^{n-1}$ partitions are distinct. Every restricted
growth word of positive length $n$ is a word over $[n]$, giving the upper bound.
The numerical values follow by the recurrence, using only integer addition
and multiplication.
\end{proof}

If $q_1,\ldots,q_n$ are distinct primes, partitions of $[n]$ correspond to
unordered factorizations of $q_1\cdots q_n$ into factors exceeding one: a block
is sent to the product of the primes it contains. In the inverse algorithm,
test divisibility of every factor by each of the given primes. The elementary
prime-divisor lemma, applied finitely often, makes these supports disjoint and
exhaustive. Thus the squarefree-factorization statement also has a finite
proof. The input includes the distinct primes; no general factorization oracle
is needed.

\subsection{Recurrences and finite summation formulas}

The two familiar recurrences are unified by Spivey's finite identity
\cite{spivey}; the proof below gives the explicit bijection.

\begin{theorem}[Bell recurrence and Spivey's identity]\label{thm:spivey}
For $n,m\geq0$, using $0^0=1$,
\begin{align}
 B_{n+1}&=\sum_{k=0}^{n}\binom nk B_k,\label{eq:bell-binomial}\\
 B_{n+m}&=\sum_{k=0}^{n}\sum_{j=0}^{m}
       \binom nk S(m,j)j^{n-k}B_k.\label{eq:spivey}
\end{align}
\end{theorem}
\begin{proof}
For the first identity, delete the entire block containing a distinguished
element of an $(n+1)$-element set. If $k$ elements remain, their set can be
chosen in $\binom nk$ ways and partitioned in $B_k$ ways. Putting the
distinguished element together with the complement reverses the operation.

For the second identity, first restrict a partition to a fixed $m$-element
subset. If this restriction has $j$ blocks, there are $S(m,j)$ choices for it.
Among the other $n$ elements choose the $k$ elements whose blocks do not meet
the first subset. Partition those elements in $B_k$ ways. Each of the remaining
$n-k$ elements chooses one of the $j$ old blocks. These choices are independent,
and give $j^{n-k}$ possibilities. Restricting a completed partition recovers
all the choices uniquely. The convention $0^0=1$ handles empty subsets.
\end{proof}

Here and below a combinatorial count is an equality between bounded sums of
Boolean characteristic functions. Each displayed inverse map gives a
primitive recursive permutation of the relevant finite codes; equality of
their sums follows by finite reindexing.

\begin{proposition}[Bell triangle]\label{prop:bell-triangle}
Define $T(0,0)=1$ and, for $0\leq k\leq n$,
\[
  T(n,0)=B_n,\qquad
  T(n,k)=T(n,k-1)+T(n-1,k-1)\quad(k>0).
\]
Then
\[
  T(n,k)=\sum_{i=0}^{k}\binom ki B_{n-i},\qquad
  T(n,n)=B_{n+1}.
\]
In particular the next row begins with the preceding row's last entry.
\end{proposition}
\begin{proof}
Pascal's identity shows that the displayed sum satisfies the defining
addition rule, including its initial column. At $k=n$, replace $i$ by $n-i$
and use \eqref{eq:bell-binomial}. These verifications use induction through a
finite triangular table.
\end{proof}

For an indeterminate $X$, put $(X)_0=1$ and
$(X)_k=X(X-1)\cdots(X-k+1)$. Comparing the recurrence in
Proposition~\ref{prop:stirling} with
$X(X)_k=(X)_{k+1}+k(X)_k$ proves the polynomial identity
\begin{equation}\label{eq:power-falling}
 X^n=\sum_{k=0}^{n}S(n,k)(X)_k.
\end{equation}
All polynomial identities here are statements about finite coefficient
lists, not assertions requiring a theory of arbitrary functions.

\begin{theorem}[Inclusion--exclusion and finite Dobinski forms]
\label{thm:finite-dobinski}
For all $n,k\geq0$,
\begin{equation}\label{eq:stirling-ie}
 k!S(n,k)=\sum_{i=0}^{k}(-1)^{k-i}\binom ki i^n.
\end{equation}
Consequently
\begin{align}
 B_n&=\sum_{k=0}^{n}\frac1{k!}
        \sum_{i=0}^{k}(-1)^{k-i}\binom ki i^n,\label{eq:bell-ie}\\
 B_n&=\sum_{i=0}^{n}\frac{i^n}{i!}
        \sum_{j=0}^{n-i}\frac{(-1)^j}{j!}.\label{eq:bell-finite-dobinski}
\end{align}
If $D_t$ is the number of fixed-point-free permutations of $[t]$, then
\begin{equation}\label{eq:bell-derangements}
  B_n=\sum_{i=0}^{n}\frac{i^nD_{n-i}}{i!(n-i)!}.
\end{equation}
\end{theorem}
\begin{proof}
There are $k!S(n,k)$ surjections from $[n]$ to $[k]$: label the blocks of a
partition bijectively by the target labels. For a fixed function with image
$I$, its contribution to the right side of \eqref{eq:stirling-ie} is
\[
 \sum_{I\subseteq A\subseteq[k]}(-1)^{k-|A|}
       =(1-1)^{k-|I|}.
\]
This is one exactly when the function is onto and zero otherwise. Interchanging
the two finite sums proves \eqref{eq:stirling-ie}; summing over $k$ proves
\eqref{eq:bell-ie}. Put $j=k-i$ and use
$\binom ki/k!=1/(i!j!)$ to obtain
\eqref{eq:bell-finite-dobinski}.

The same inclusion--exclusion calculation for forbidden fixed points gives
$D_t=t!\sum_{j=0}^{t}(-1)^j/j!$: imposing a specified set of $j$ fixed points
leaves $(t-j)!$ permutations. Substitution gives
\eqref{eq:bell-derangements}. These rational formulas are finite arithmetic
equalities; multiplication by a product of the relevant factorials removes
all division.
\end{proof}

One may also group partitions by their block sizes. If $m_j$ blocks have size
$j$, then the labels can be assigned to ordered slots in $n!$ ways. Permuting
within each block and permuting equal-sized blocks gives exactly
$\prod_{j=1}^{n}(j!)^{m_j}m_j!$ assignments with the same partition. Hence
\begin{equation}\label{eq:block-profile}
 B_n=\sum_{\substack{m_1,\ldots,m_n\leq n\\\sum_{j=1}^n jm_j=n}}
       \frac{n!}{\prod_{j=1}^{n}(j!)^{m_j}m_j!}.
\end{equation}
The quotient is exact, justified by this explicit finite grouping. More
generally, weighting a block of size $j$ by $x_j$ defines the complete Bell
polynomial
\[
 C_n(x_1,\ldots,x_n)=\sum_{\pi}\prod_{A\in\pi}x_{|A|}.
\]
Its coefficient sum is $C_n(1,\ldots,1)=B_n$.
The one-variable Touchard polynomial
$\Phi_n(t)=\sum_kS(n,k)t^k$ similarly satisfies
\[
 \Phi_n(1)=B_n,\qquad
 \Phi_{n+1}(t)=t\bigl(\Phi_n(t)+\Phi_n'(t)\bigr),
\]
by the Stirling recurrence. The derivative is the formal derivative of a
finite polynomial.

\subsection{A finite moment functional and the shift identities}
\label{subsec:finite-functional}

Because the falling-factorial polynomials are monic of successively increasing
degrees, every polynomial $f\in\Q[X]$ has a unique finite expansion
$f=\sum_jc_j(X)_j$. Define
\[
 \mathcal L(f)=\sum_jc_j.
\]
Triangular coefficient elimination computes $\mathcal L$. Identity
\eqref{eq:power-falling} says $\mathcal L(X^n)=B_n$, while by definition
$\mathcal L((X)_j)=1$.

\begin{lemma}[Finite shift identity]\label{lem:finite-shift}
For every polynomial $f$ and every $k\geq0$,
\begin{equation}\label{eq:finite-shift}
 \mathcal L\bigl((X)_k f(X)\bigr)
          =\mathcal L\bigl(f(X+k)\bigr).
\end{equation}
\end{lemma}
\begin{proof}
First let $g(Y)=(Y)_j$. Then
$(X)_k g(X-k)=(X)_{k+j}$, so applying $\mathcal L$ gives one on both sides of
$\mathcal L((X)_k g(X-k))=\mathcal L(g(X))$. Expand an arbitrary $g$ in the
falling-factorial basis and use linearity. Finally set $g(Y)=f(Y+k)$.
\end{proof}

This single finite lemma proves the additional summation identities on the
Wikipedia page, attributed there to Komatsu--Pita-Ruiz~\cite{komatsu-pita}.
For rational $a,b$ and natural $k,n$,
\begin{equation}\label{eq:general-umbral-shift}
 \mathcal L\bigl((aX+b)^n\bigr)
 =\mathcal L\bigl((X)_k(aX+b-ak)^n\bigr).
\end{equation}
Define the signed first-kind Stirling coefficients by
$(X)_k=\sum_i s(k,i)X^i$, equivalently
$s(k,i)=(-1)^{k-i}\genfrac{[}{]}{0pt}{}{k}{i}$. Expanding both sides gives
\begin{align}\label{eq:general-stirling-shift}
 &\sum_{j=0}^{n}\binom nj a^j b^{n-j}B_j\notag\\
 &\quad=\sum_{i=0}^{k}s(k,i)
       \sum_{j=0}^{n}\binom nj a^j(b-ak)^{n-j}B_{j+i}.
\end{align}
In particular $k=1$ gives the displayed one-step identity of the source,
and $a=1,b=k$ gives
\[
 \sum_{j=0}^{n}\binom nj k^{n-j}B_j
       =\sum_{i=0}^{k}s(k,i)B_{n+i}.
\]
Using \eqref{eq:finite-shift} with $k=1$ and
$f(X)=X^n(X-1)^r$ also gives
\[
 \mathcal L\bigl(X^{n+1}(X-1)^r\bigr)
       =\mathcal L\bigl(X^r(X+1)^n\bigr).
\]
On expanding, this is
\begin{equation}\label{eq:general-pascal-inversion}
 \sum_{j=0}^{n}\binom nj B_{r+j}
   =\sum_{i=0}^{r}(-1)^{r-i}\binom ri B_{n+i+1}.
\end{equation}
Its specialization $n=0$ is Pascal inversion for the Bell recurrence.
No infinite moment representation was needed for these identities.

\subsection{The Hankel determinant and strict log-convexity}

\begin{theorem}[Hankel determinant]\label{thm:hankel}
For every $N\geq0$,
\begin{equation}\label{eq:hankel}
 \det\bigl(B_{i+j}\bigr)_{0\leq i,j\leq N}
       =\prod_{j=0}^{N}j!.
\end{equation}
\end{theorem}
\begin{proof}
Define monic Charlier polynomials by the finite formula
\[
 P_r(X)=\sum_{j=0}^{r}(-1)^{r-j}\binom rj(X)_j.
\]
Let $\Delta f(X)=f(X+1)-f(X)$. The finite shift identity and the binomial
formula for an iterated difference yield
\[
 \mathcal L(P_r f)
   =\sum_{j=0}^{r}(-1)^{r-j}\binom rj\mathcal L(f(X+j))
   =\mathcal L(\Delta^r f).
\]
If $\deg f<r$, its $r$th difference is zero. If $f$ is monic of degree $r$,
the $r$th difference is the constant $r!$. Applying this to the $P_s$, and
using symmetry when $r<s$, gives
\[
  \mathcal L(P_rP_s)=\begin{cases}r!&r=s,\\0&r\ne s.\end{cases}
\]
The coefficient matrix $U$ expressing $P_0,\ldots,P_N$ in the monomial basis
is lower triangular with diagonal entries one. Thus, by multiplication of
finite matrices,
\[
 U\,(B_{i+j})_{i,j=0}^{N}\,U^{\mathsf T}
       =\operatorname{diag}(0!,1!,\ldots,N!).
\]
Taking determinants, defined by their finite permutation sums, proves
\eqref{eq:hankel}, since $\det U=1$.
\end{proof}

\begin{corollary}[Strict log-convexity]\label{cor:strict-logconvex}
For $n\geq1$,
\begin{equation}\label{eq:strict-logconvex}
 B_n^2<B_{n-1}B_{n+1}.
\end{equation}
Consequently the quotients $B_{n+1}/B_n$ increase, and
$B_mB_n\leq B_{m+n}$ for $m,n\geq0$.
\end{corollary}
\begin{proof}
If $f=\sum_jc_jP_j$, the orthogonality calculation gives
$\mathcal L(f^2)=\sum_jj!c_j^2$, positive for nonzero $f$.
For two independent polynomials $f,g$, write their coefficient vectors in this
basis as $c,d$. Then the entirely finite identity
\[
 \mathcal L(f^2)\mathcal L(g^2)-\mathcal L(fg)^2
       =\sum_{i<j}i!j!(c_id_j-c_jd_i)^2>0
\]
follows by expansion. At least one minor is nonzero, since otherwise the
vectors would be proportional. The same conclusion holds for the bilinear
form $\mathcal L(Xfg)$, because \eqref{eq:finite-shift} identifies it with
$\mathcal L(f(X+1)g(X+1))$, and substitution $X\mapsto X+1$ preserves
independence.

Write $n-1=2q+\epsilon$ with $\epsilon\in\{0,1\}$. Take $f=X^q$ and
$g=X^{q+1}$, using the ordinary bilinear form for $\epsilon=0$ and the shifted
one for $\epsilon=1$. Their Gram determinant is precisely
$B_{n-1}B_{n+1}-B_n^2$, proving strict positivity. Finally multiply the
inequalities
$B_{n+j+1}/B_{n+j}\geq B_{j+1}/B_j$ for $0\leq j<m$; the products telescope.
All denominators are positive integers and can be cleared.
\end{proof}

\subsection{Normalized log-concavity, with a finite proof}
\label{subsec:logconcavity}

The complementary inequality also admits a finite proof. We give the
coefficient argument underlying the Bender--Canfield method
\cite{bender-canfield}, specialized enough that its arithmetic content is
visible.

\begin{lemma}[A finite coefficient inequality]\label{lem:coefficient-logconcave}
Let $x_0=1,x_1,x_2,\ldots$ be nonnegative rational numbers such that
\[
 x_jx_{k-1}\geq x_{j-1}x_k\qquad(1\leq j<k).
\]
Let $a_0=1$, put $a_t=0$ for $t<0$, and define
\[
 r a_r=\sum_{j=1}^{r}x_j a_{r-j}\qquad(r\geq1).
\]
Then, for $1\leq m\leq n$,
\begin{equation}\label{eq:general-logconcavity}
 a_m a_n\geq a_{m-1}a_{n+1}.
\end{equation}
Only $x_0,\ldots,x_{n+1}$ and $a_0,\ldots,a_{n+1}$ occur in each assertion.
\end{lemma}
\begin{proof}
The recurrence, with its $x_0a_r$ term restored, says
$\sum_{j=1}^{r+1}x_{j-1}a_{r+1-j}=(r+1)a_r$.
Expand the product of this identity and the defining recurrence, subtract
the two resulting products, and pair indices $j,k$. One obtains the finite
identity
\begin{align}\label{eq:paired-logconcavity}
 &m(n+1)(a_m a_n-a_{m-1}a_{n+1})\notag\\
 &\quad=\sum_{1\leq j<k\leq n+1}
 (x_jx_{k-1}-x_{j-1}x_k)
 (a_{m-j}a_{n+1-k}-a_{m-k}a_{n+1-j}).
\end{align}
Indeed, before pairing, the summand is
$(x_jx_{k-1}-x_{j-1}x_k)a_{m-j}a_{n+1-k}$; summands with $j>m$ vanish by the
negative-index convention, and exchanging $j,k$ negates the first factor.

Induct on $m$, retaining all $n\geq m$ as free parameters. If $k>m$, the
second subtracted product in \eqref{eq:paired-logconcavity} is zero, so its
bracket is nonnegative. Suppose $k\leq m$, and put
\[
 A=m-j,\quad C=m-k,\quad D=n+1-j,\quad E=n+1-k.
\]
Then $C\leq\min(A,E)$, $\max(A,E)\leq D$, and $A+E=C+D$.
Starting with the pair $(\min(A,E),\max(A,E))$, move one unit at a time
from its smaller index to its larger one until the pair becomes $(C,D)$.
At each step the smaller index is less than $m$, so the induction hypothesis
says the product of the two $a$'s cannot increase. Thus
$a_Aa_E\geq a_Ca_D$. The case $m=1$ has only the already-settled
$k>m$ terms. Every summand in \eqref{eq:paired-logconcavity} is therefore
nonnegative. Since $m(n+1)>0$, this proves the lemma.

For a literal quantifier-free induction implementation, fix a final bound
$N$ and retain the Boolean conjunction of these inequalities for all
$1\leq m\leq n\leq N$ reached so far. The above steps never require an index
exceeding $N+1$. Clearing the product of denominators makes the invariant an
integer inequality test on a finite table.
\end{proof}

\begin{theorem}[Normalized log-concavity]\label{thm:normalized-logconcavity}
For every $n\geq1$,
\begin{equation}\label{eq:normalized-logconcavity}
 n B_{n-1}B_{n+1}\leq(n+1)B_n^2.
\end{equation}
Equivalently, $B_n/n!$ is log-concave. Moreover,
\begin{equation}\label{eq:bell-product-bounds}
 B_mB_n\leq B_{m+n}\leq\binom{m+n}{m}B_mB_n.
\end{equation}
\end{theorem}
\begin{proof}
Set $a_n=B_n/n!$, $x_0=1$, and $x_j=1/(j-1)!$ for $j\geq1$.
Dividing \eqref{eq:bell-binomial} by $n!$ gives
$(n+1)a_{n+1}=\sum_{j=1}^{n+1}x_j a_{n+1-j}$.
The successive ratios $x_j/x_{j-1}$ are
$1,1,1/2,1/3,\ldots$, a decreasing sequence, so the hypothesis of
Lemma~\ref{lem:coefficient-logconcave} holds. Its case $m=n$ is exactly
\eqref{eq:normalized-logconcavity} after factorials are cleared.

For the upper product bound, normalized log-concavity makes $a_{j+1}/a_j$
decrease. Hence
\[
 \frac{a_{n+m}}{a_n}
 =\prod_{j=0}^{m-1}\frac{a_{n+j+1}}{a_{n+j}}
 \leq\prod_{j=0}^{m-1}\frac{a_{j+1}}{a_j}=a_m.
\]
Clearing factorials gives the upper inequality in
\eqref{eq:bell-product-bounds}; the lower one was proved in
Corollary~\ref{cor:strict-logconvex}. The empty product handles $m=0$.
\end{proof}

\subsection{Permutation interpretations and card shuffles}

A generalized pattern $1$--$23$ means indices $i<j<j+1$ with
$\pi_i<\pi_j<\pi_{j+1}$; $1$--$32$ instead means
$\pi_i<\pi_{j+1}<\pi_j$. These definitions are bounded tests on a permutation.

\begin{proposition}[Eight avoidance classes]\label{prop:pattern-avoidance}
For each of
\[
 1\text{--}23,\quad32\text{--}1,\quad3\text{--}21,\quad12\text{--}3,
 \qquad
 1\text{--}32,\quad23\text{--}1,\quad3\text{--}12,\quad21\text{--}3,
\]
the number of avoiding permutations of $[n]$ is $B_n$.
\end{proposition}
\begin{proof}
Given a partition, order its blocks by decreasing minima. For the first
class, write each block with its minimum first and its remaining elements in
decreasing order, then erase the block separators. An ascent can occur only
immediately after a left-to-right minimum, so it cannot be preceded by a
smaller element. The result avoids $1$--$23$. Conversely, cut such an avoiding
permutation immediately before each left-to-right minimum. Within a segment,
every adjacent pair after its first entry must decrease: otherwise the first
entry would complete a forbidden pattern. The segments therefore recover
exactly the original partition and its prescribed order.

For the class $1$--$32$, use the same block order but write every block in
increasing order. A descent can then end only at a left-to-right minimum.
This is exactly the condition excluding $1$--$32$: a descent ending at an
element that is not a new minimum has an earlier smaller element and gives
the forbidden pattern. Cutting at the minima again supplies the inverse.
Reversing the positions and complementing every value $v$ to $n-1-v$ are
explicit involutions on permutations. Their four combinations take the first
pattern to the first four listed patterns, and the second to the last four.
\end{proof}

\begin{proposition}[Extending every $321$ to a $3241$]
\label{prop:barred-pattern}
There are $B_n$ permutations of $[n]$ in which every triple of positions
$i<j<k$ with $\pi_i>\pi_j>\pi_k$ has a position $j<\ell<k$ satisfying
$\pi_\ell>\pi_i$.
\end{proposition}
\begin{proof}
Cut a permutation immediately before each left-to-right maximum. Call these
its record segments; each begins with its largest element. The stated
property is equivalent to the tail of every record segment being increasing.
For necessity, a decreasing pair in such a tail, together with the leading
record, is a $321$ that cannot be extended: no larger value occurs before
the end of the segment. Conversely, consider any $321$. Its middle element
is not a left-to-right maximum. If its last element lay in the same segment,
the increasing-tail condition would be violated. Otherwise a new record lies
between the middle and last elements, and this record exceeds the first
element of the triple, supplying the required fourth position.

Given a set partition, order blocks by increasing maxima and write each block
with its maximum first and the remaining elements in increasing order.
Erasing separators produces exactly a permutation with the just-described
record segments. Cutting at the records is its inverse. This is the required
finite bijection.
\end{proof}

\begin{proposition}[Return shuffles]\label{prop:return-shuffles}
For a deck of $n\geq1$ distinct cards, there are $n^n$ sequences of $n$
operations that remove the top card and reinsert it in any of the $n$
positions. Exactly $B_n$ return the deck to its original order. Thus the
return probability under independent uniform choices is $B_n/n^n$.
\end{proposition}
\begin{proof}
Reverse a path of operations. An inverse operation chooses a card and moves
it to the top. Label the original order $1,\ldots,n$. After a word of such
choices, the chosen cards occur first, in reverse order of their last
occurrences in the word; untouched cards follow in their original order.
This description follows by induction on the length of the word.

To return to $1,\ldots,n$, the chosen labels must therefore be a prefix
$\{1,\ldots,k\}$, and their last occurrences must be ordered from latest to
earliest as $1,\ldots,k$. Partition the $n$ action times according to equal
chosen labels. Conversely, given any partition of those times, label its
blocks $1,\ldots,k$ in decreasing order of their maxima. The resulting choice
word returns the deck to its original order. The constructions are inverse;
reversing operations now proves the assertion for top-card reinsertion.
\end{proof}

The permutation bijection gives $B_n\leq n!$. A second useful comparison is
$n^n\leq n!B_n$: partition functions $[n]\to[n]$ by their fibers and use
\eqref{eq:power-falling} at $X=n$, noting $(n)_k\leq n!$.
For $n\geq3$ this is strict, since the $k=1$ term has
$(n)_1=n<n!$. Consequently the return probability exceeds $1/n!$ for
$n\geq3$; equality holds for $n=1,2$.

\paragraph{Arithmetic audit.}
The largest finite objects used in this section are tables, words,
permutations, and coefficient lists with explicitly bounded lengths.
Their enumerations have primitive recursive bounds, such as $(n+1)^n$, $n!$,
and products of factorials. Every appeal to induction can retain a finite
table of the preceding assertions as one quantifier-free invariant. The
positive-denominator convention turns all rational inequalities into integer
inequalities. Thus the arguments use no unrestricted induction, choice,
excluded middle for undecidable formulas, or existence principle for infinite
sets.

% Bibliography metadata:
% bender-canfield: E. A. Bender and E. R. Canfield,
% Log-Concavity and Related Properties of the Cycle Index Polynomials,
% J. Combin. Theory Ser. A 74 (1996), 57--70,
% https://mathweb.ucsd.edu/~ebender/reprints/93.pdf
% Additional finite-sum statements match T. Komatsu and C. Pita-Ruiz,
% Some formulas for Bell numbers, Filomat 32(11) (2018), 3881--3889,
% https://doi.org/10.2298/FIL1811881K . Their proofs above follow directly
% from the finite shift functional, without using the article as a premise.

\section{Finite identities behind the analytic representations}\label{sec:series}

Throughout, $0^0=1$, $B_0=1$, and a sum over an empty range is zero.  The
Stirling numbers and Bell numbers are the primitive recursive integers defined
in Section~\ref{sec:finite}.  We first prove exact truncation identities; the
infinite notation is then a consequence with an explicit error, not a premise.

\subsection{Dobi\'nski's formula with a uniform arithmetic error}

Write $(x)_j=x(x-1)\cdots(x-j+1)$ and $(x)_0=1$.  The polynomial identity
\[
 x^n=\sum_{j=0}^n S(n,j)(x)_j
\]
follows either from the Stirling recurrence or by classifying functions
$[n]\to[x]$ according to their nonempty fibers.  Substituting the integers
$k=0,\ldots,M$ and interchanging two \emph{finite} sums gives, for $M\ge n$,
\begin{equation}\label{eq:finite-dobinski}
 A_{n,M}:=\sum_{k=0}^M\frac{k^n}{k!}
 =\sum_{j=0}^n S(n,j) E_{M-j}(1).
\end{equation}
The cancellation $(k)_j/k!=1/(k-j)!$ is used only when $k\ge j$; for $k<j$
the falling factorial vanishes.

\begin{theorem}[Dobi\'nski, with a finite certificate]\label{thm:dobinski}
For $M\ge n$,
\begin{equation}\label{eq:dobinski-error}
 0\le B_n-e^{-1}A_{n,M}\le
 \frac{2B_n}{(M-n+1)!}.
\end{equation}
Consequently
\[
 B_n=e^{-1}\sum_{k=0}^{\infty}\frac{k^n}{k!}.
\]
For absolute error at most $2^{-s}$, it is sufficient to take
$M=n+s+(n+1)^2+2$, and to approximate $e^{-1}$ to the additional precision
needed for multiplication by the explicit rational $A_{n,M}$.
\end{theorem}
\begin{proof}
Subtract \eqref{eq:finite-dobinski} from
$eB_n=e\sum_{j=0}^nS(n,j)$.  Each difference $e-E_{M-j}(1)$ is nonnegative
and at most $2/(M-j+1)!\le2/(M-n+1)!$.  Sum the inequalities and divide by
$e\ge1$.  Finally, $B_n\le(n+1)^{n+1}\le2^{(n+1)^2}$ and
$d!\ge2^{d-1}$ for $d\ge1$, both by finite induction.  These inequalities
verify the claimed truncation rule.  If the numerical value of $A_{n,M}$ is
large, choose $t=s+\lceil A_{n,M}\rceil+3$ for the reciprocal's precision;
then $A_{n,M}2^{-t}\le2^{-s-2}$.  Every operation in this rule is primitive
recursive.
\end{proof}

\begin{corollary}[The Poisson interpretation and its polynomial extension]
The weights $w_k=e^{-1}/k!$ have total mass $1$ and their $n$th moment is $B_n$.
For a nonnegative rational parameter $u$, define
$\Phi_n(u)=\sum_{j=0}^n S(n,j)u^j$.  Then
\begin{equation}\label{eq:poisson-touchard}
 e^{-u}\sum_{k=0}^{\infty}\frac{k^n u^k}{k!}=\Phi_n(u),
 \qquad \Phi_{n+1}(u)=u\bigl(\Phi_n(u)+\Phi_n'(u)\bigr).
\end{equation}
\end{corollary}
\begin{proof}
For finite $M\ge n$ the weighted version of
\eqref{eq:finite-dobinski} is
$\sum_{k\le M}k^nu^k/k!=\sum_{j\le n}S(n,j)u^j E_{M-j}(u)$.
Apply \eqref{eq:exp-tail}.  The derivative identity follows by substituting
$S(n+1,j)=jS(n,j)+S(n,j-1)$ and collecting polynomial coefficients.  There is
no construction of a sigma-algebra or arbitrary random variable here:
``moment'' denotes this specified effective weighted sum.
\end{proof}

\subsection{A finite ceiling formula}

\begin{theorem}[Comtet's truncation]\label{thm:comtet}
For every integer $n\ge1$,
\[
 B_n=\left\lceil e^{-1}\sum_{k=1}^{2n}\frac{k^n}{k!}\right\rceil.
\]
\end{theorem}
\begin{proof}
Set $t_k=k^n/k!$.  For $k\ge2n+1$,
\[
 \frac{t_{k+1}}{t_k}
 =\frac{(1+1/k)^n}{k+1}
 \le\frac{e^{n/k}}{k+1}
 <\frac2{k+1}\le\frac12.
\]
Here $e^{1/2}<2$, for example by squaring and using $e<3<4$.
Let $a_n=(2n+1)^n/(2n+1)!$.  Then $a_1=1/2$ and
\[
 \frac{a_{n+1}}{a_n}
 =\frac{(1+2/(2n+1))^n}{2n+2}
 <\frac3{2n+2}\le\frac34.
\]
Thus $a_n\le1/2$, and the omitted positive tail is less than $2a_n\le1$.
By Theorem~\ref{thm:dobinski}, the truncated expression lies strictly between
$B_n-1$ and $B_n$.  This is precisely the ceiling assertion.  The constants
$e<3$ and $(1+x)^m\le e^{mx}$ follow from the factorial series and the binomial
expansion.  This proof does not assume that the ceiling of an arbitrary real
name is computable: it proves a specific interval with an integer endpoint.
\end{proof}

\subsection{Exponential generating functions and complete Bell polynomials}

A formal power series identity below means equality of every finite truncated
coefficient list.  In $\Q[z]/(z^{N+1})$, substitution into an exponential is
finite whenever the substituted polynomial has constant coefficient zero.

\begin{theorem}[The exponential generating function]\label{thm:egf}
As a formal series, and also as an effective complex analytic identity,
\begin{equation}\label{eq:egf}
 \sum_{n=0}^{\infty}\frac{B_n}{n!}z^n=\exp(e^z-1).
\end{equation}
In particular, for every real $r>0$ and every $n\ge0$,
\begin{equation}\label{eq:coefficient-bound}
 B_n\le n!\,r^{-n}\exp(e^r-1).
\end{equation}
\end{theorem}
\begin{proof}
The coefficient of $z^n$ in $(e^z-1)^k/k!$ counts ordered lists of $k$
nonempty labeled blocks, divided by their $k!$ possible orders; it is
$S(n,k)/n!$.  In detail, choose block sizes
$d_1+\cdots+d_k=n$ with $d_i\ge1$.  The multinomial count is
$n!/(d_1!\cdots d_k!)$; dividing by $n!k!$ gives exactly the coefficient
obtained by multiplying the $k$ series.  All $k>n$ contribute zero.
Summing $k=0,\ldots,n$ proves the formal identity.

There is also a direct recurrence check: if $F(z)=\sum B_nz^n/n!$, the
binomial recurrence gives $F'=e^zF$ and $F(0)=1$.  The formal exponential on
the right of \eqref{eq:egf} has the same equations.  Coefficient induction
therefore gives the same result without formal division by an unknown series.

For $r>0$, every finite partial sum of $F(r)$ is bounded above by
$\sum_{k\ge0}(e^r-1)^k/k!=\exp(e^r-1)$: expand the finite coefficients just
computed, and only add nonnegative terms.  The single $n$th term yields
\eqref{eq:coefficient-bound}.  For $|z|\le A$ and $R=2(A+1)$, the coefficient
bound gives
\begin{equation}\label{eq:egf-tail}
 \sum_{n>N}\frac{B_n}{n!}|z|^n
 \le 2\exp(e^R-1)\,2^{-(N+1)}.
\end{equation}
An explicit integer upper bound for the large constant is available by
$e<3$: for integral $R$, take
$\exp(e^R-1)<3^{3^R}$.  Hence choosing
$N=s+2\cdot3^R+4$ suffices.  For rational $R$, first increase it to an integer.
The finite coefficient identity, the exponential tails, and
\eqref{eq:egf-tail} now justify the analytic identity.  They also justify all
termwise differentiations on smaller bounded disks, by inserting the extra
polynomial factors $n(n-1)\cdots$ in the same geometric tail argument.
\end{proof}

The complete Bell polynomial can be defined without analytic vocabulary by
\[
 C_n(x_1,\ldots,x_n)
 =\sum_{\pi\text{ a partition of }[n]}
       \prod_{C\in\pi}x_{|C|}.
\]
Grouping by the number $m_j$ of size-$j$ blocks proves the finite formula
\begin{equation}\label{eq:block-multiplicity}
 C_n(x_1,\ldots,x_n)
 =\sum_{\substack{m_1,\ldots,m_n\ge0\\\sum j m_j=n}}
   \frac{n!}{\prod_{j=1}^n m_j!(j!)^{m_j}}
   \prod_{j=1}^n x_j^{m_j}.
\end{equation}
Indeed, order all $n$ labels, cut them into the indicated blocks, and quotient
by the $j!$ orders within each block and the $m_j!$ orders of equal-sized blocks.
Thus $C_n(1,\ldots,1)=B_n$, so the sum of its coefficients is $B_n$.
Expanding a finite exponential also proves
\[
 \sum_{n\ge0}C_n(x_1,\ldots,x_n)\frac{z^n}{n!}
 =\exp\left(\sum_{j\ge1}x_j\frac{z^j}{j!}\right)
\]
coefficientwise.  In the usual moment--cumulant language, all the cumulants of
the specified Poisson weights are $1$.  This is exactly the preceding formal
identity specialized to $x_j=1$.

\subsection{Hypergeometric notation, without extra functions}

For $n\ge1$, the hypergeometric statement on MathWorld is
\[
 B_n=e^{-1}\,{}_{n-1}F_{n-1}
     (\underbrace{2,\ldots,2}_{n-1};
      \underbrace{1,\ldots,1}_{n-1};1).
\]
Define the expression on the right by its factorial series.  Since the rising
factorials satisfy $(2)^{\overline j}/(1)^{\overline j}=j+1$, its $j$th term is
$(j+1)^{n-1}/j!=(j+1)^n/(j+1)!$.  Changing the index in
Theorem~\ref{thm:dobinski} proves the formula with the same effective tail.
For $n=1$ this includes ${}_0F_0(1)=e$.

\subsection{The ordinary series and its negative-axis substitute}

The ordinary generating series $\sum B_nx^n$ is formal.  It is not a
convergent series away from $0$.  Nevertheless, on the negative real axis
there is a useful function having that series as its asymptotic expansion.

\begin{theorem}[A Stieltjes-type representation and its full remainder]
For $x\le0$, define the effective name
\[
 G(x)=e^{-1}\sum_{k=0}^{\infty}\frac1{(1-kx)k!}.
\]
It is uniformly convergent on this half-line, with tail at most
$2/(M+1)!$ after $k=M$.  For every $N\ge0$,
\begin{equation}\label{eq:ordinary-asymptotic}
 G(x)=\sum_{j=0}^N B_jx^j+\mathcal R_N(x),\qquad
 |\mathcal R_N(x)|\le B_{N+1}|x|^{N+1}.
\end{equation}
The remainder has sign $(-1)^{N+1}$ for $x<0$.
\end{theorem}
\begin{proof}
Use the exact finite identity
\[
 \frac1{1-kx}=\sum_{j=0}^N(kx)^j+
                    \frac{(kx)^{N+1}}{1-kx}.
\]
Multiply by $e^{-1}/k!$, sum first over a finite range of $k$, and apply
Theorem~\ref{thm:dobinski}.  For $x\le0$ the denominator is at least $1$,
so the absolute value of the last sum is at most $B_{N+1}|x|^{N+1}$.
The factorial tail proves existence of the name independently of the
asymptotic expansion.  Each summand in the last sum has the claimed sign.
\end{proof}

For $x<0$ put $a=-1/x>0$.  Then
$1/(1-kx)=a/(a+k)=(a)^{\overline k}/(a+1)^{\overline k}$.
Consequently the confluent hypergeometric representation is exactly
\[
 G(x)=e^{-1}\,{}_1F_1(-1/x;1-1/x;1),
\]
with the defining series and the preceding tail bound.  We do not use the
incomplete-gamma expression from MathWorld, whose written form involves
complex branch conventions.

The other convergent series displayed there also follows by finite algebra:
\begin{equation}\label{eq:ordinary-rising-series}
 G(x)=\sum_{k=0}^{\infty}\frac{(-1)^k}{(a+1)^{\overline k}}
     =\sum_{k=0}^{\infty}
       \frac{x^k}{(-x)^k((x-1)/x)^{\overline k}}
       \qquad(x<0,\ a=-1/x).
\end{equation}
Here the rising factorial with index zero is one. To prove the identity,
define the finite rational sum
\[
 H_n(a)=\sum_{j=0}^{n}(-1)^j\binom nj\frac{a}{a+j}.
\]
Pascal's identity gives
\[
 H_{n+1}(a)=H_n(a)-\frac{a}{a+1}H_n(a+1).
\]
Starting with $H_0(a)=1$, induction therefore gives
\begin{equation}\label{eq:finite-beta-sum}
 H_n(a)=\frac{n!}{(a+1)^{\overline n}}:
\end{equation}
indeed, the difference of the two proposed terms has numerator
$n!((a+n+1)-a)=(n+1)!$ over
$(a+1)\cdots(a+n+1)$. For quantifier-free induction, fix the final index
$N$ and retain the finite table of assertions for arguments
$a,a+1,\ldots,a+N-r$ at stage $r$; each next entry uses two entries already
in that table. All denominators are positive and may be cleared.

The coefficient of $z^n$ in the finite product
$e^{-z}\,{}_1F_1(a;a+1;z)$ is
\[
 \frac1{n!}\sum_{j=0}^{n}(-1)^{n-j}\binom nj\frac{a}{a+j}
   =\frac{(-1)^n}{(a+1)^{\overline n}},
\]
by \eqref{eq:finite-beta-sum}. Since $0<a/(a+j)\leq1$ and
$(a+1)^{\overline n}\geq n!$, all three series and their product have
factorial tail bounds on every fixed bounded disk. Thus the finite coefficient
calculation passes to their effective names. Setting $z=1$ proves
\eqref{eq:ordinary-rising-series}, with absolute tail at most
$2/(M+1)!$ after index $M$. This adds no special-function assumption.

To see directly that the ordinary series has radius zero, fix a positive
rational $x$ and choose an integer $k$ with $kx\ge2$.  Dobi\'nski gives
\[
 B_nx^n\ge\frac{(kx)^n}{e\,k!}\ge\frac{2^n}{3k!}.
\]
An explicit $n$ makes this greater than any prescribed rational number.
The terms therefore do not tend to zero.  The same absolute-value argument
applies at nonzero complex $x$ with an explicit positive rational lower bound
for $|x|$.  This is also an elementary proof that Bell numbers grow faster
than every fixed exponential, independent of the sharper growth theorems.

\subsection{Fourier coefficient integrals and Ces\`aro's formula}

\begin{theorem}[Circular coefficient extraction]\label{thm:fourier-coefficient}
For $r>0$ and $n\ge0$,
\begin{equation}\label{eq:circular-extraction}
 B_n=\frac{n!}{2\pi r^n}
  \int_{-\pi}^{\pi}\exp(e^{r e^{i\theta}}-1)e^{-in\theta}\,d\theta.
\end{equation}
Equivalently, on the positively oriented circle $|z|=r$,
\[
 B_n=\frac{n!}{2\pi i e}\int_{|z|=r}
                    \frac{\exp(e^z)}{z^{n+1}}\,dz.
\]
For $n\ge1$ the real integral representation is
\begin{align}\label{eq:cesaro}
 B_n=\frac{2n!}{\pi e}\int_0^\pi
  &\exp\!\bigl(e^{\cos\theta}\cos(\sin\theta)\bigr)\notag\\[-3pt]
  &\quad\cdot\sin\!\bigl(e^{\cos\theta}\sin(\sin\theta)\bigr)
    \sin(n\theta)\,d\theta.
\end{align}
\end{theorem}
\begin{proof}
Insert the uniformly convergent series \eqref{eq:egf} on the circle.  The
finite orthogonality identity
$(2\pi)^{-1}\int_{-\pi}^{\pi}e^{i(m-n)\theta}\,d\theta=[m=n]$
extracts the coefficient, and \eqref{eq:egf-tail} bounds the discarded part.
Parametrizing the circle by $z=r e^{i\theta}$ proves the contour notation.
No general Cauchy theorem is required.

At $r=1$ the imaginary part of $\exp(e^{e^{i\theta}})$ is
$e\sum_{m\ge1}B_m\sin(m\theta)/m!$.  Multiplication by $\sin(n\theta)$
and the identity
$\int_0^\pi\sin(m\theta)\sin(n\theta)\,d\theta=(\pi/2)[m=n]$
give the result.  Finally expand the two nested complex exponentials into
their sine and cosine parts to obtain \eqref{eq:cesaro}.  This agrees with
the corrected formula in Callan~\cite{callan-cesaro}; the factor $n!$ is essential.
All integral identities here are identities of the effective Riemann-sum
names in Section~\ref{subsec:names}.
\end{proof}

\section{Congruences and constructive periods}
\label{sec:congruences}

All congruences in this section are decidable integer assertions.
Polynomial calculations modulo $m$ mean coefficient-list calculations with
entries in $\{0,\ldots,m-1\}$. For a polynomial $F(X)=\sum_jc_jX^j$, use the
finite shorthand
\[
 (F(E)B)_n=\sum_jc_jB_{n+j}.
\]
The letter $E$ denotes a shift in the subscript, not an additional object
quantified over in the arithmetic language.

\subsection{Touchard's congruence}

\begin{theorem}[Touchard]\label{thm:touchard}
If $p$ is prime and $n\geq0$, then
\begin{equation}\label{eq:touchard}
 B_{n+p}\equiv B_{n+1}+B_n\pmod p.
\end{equation}
In particular $B_p\equiv2\pmod p$.
\end{theorem}
\begin{proof}
Let a cyclic permutation rotate $p$ distinguished elements of an
$(n+p)$-element set and fix the other $n$ elements. It acts, by a bounded
relabel-and-normalize algorithm, on the finite codes of partitions. Each
orbit has size one or $p$: the least positive return time divides the prime
$p$. One can find that return time by testing the $p$ possibilities. Choosing
the least code of each orbit therefore shows that, modulo $p$, only fixed
partitions contribute to their total number.

We classify the fixed partitions. A block meeting a fixed element must be
fixed as a block, since its rotate contains the same fixed element. If that
block meets the rotating set, it consequently contains the entire rotating
set. More generally, the restriction of a fixed partition to the $p$
rotating elements is a translation-invariant equivalence relation on
$\mathbb Z/p\mathbb Z$. The class of zero is closed under differences:
if $a$ and $b$ are equivalent to zero, translating $a\sim b$ by $-b$ gives
$a-b\sim0$. Hence it is a subgroup. A subgroup containing a nonzero element
contains all residues, because its first $p$ multiples are distinct;
otherwise it is $\{0\}$. Thus the rotating elements are either all in one
block or all in singleton blocks.

In the first case, collapse their common block contribution to one
distinguished element; this gives exactly $B_{n+1}$ partitions, allowing that
element to join any fixed elements. In the singleton case no rotating element
can share a block with a fixed element, by the preceding observation, leaving
exactly $B_n$ choices. The classes are disjoint. This proves
\eqref{eq:touchard}; setting $n=0$ gives the last assertion.
\end{proof}

There is also a short algebraic proof that is particularly convenient for
formalization. Lemma~\ref{lem:finite-shift} gives
\[
 \mathcal L\bigl(((X)_p-1)X^n\bigr)
  =\mathcal L((X+p)^n-X^n)\equiv0\pmod p.
\]
In $\mathbb F_p[X]$ one has
\begin{equation}\label{eq:finite-field-product}
 (X)_p=\prod_{a=0}^{p-1}(X-a)=X^p-X.
\end{equation}
To verify this last equality, the binomial theorem proves
$(a+1)^p\equiv a^p+1\pmod p$; induction through the residues gives
$a^p=a$ for all $a\in\mathbb F_p$. Both monic polynomials in
\eqref{eq:finite-field-product} have degree $p$ and these $p$ roots. Repeated
division by $X-a$ is legitimate because distinct residues have nonzero,
invertible differences, so their difference of degree less than $p$ is zero.
This proves \eqref{eq:touchard} by finite coefficient manipulation.

\subsection{Prime-power shifts and the prime period}

\begin{theorem}[Generalized Touchard congruence]\label{thm:touchard-power}
For prime $p$ and $r,n\geq0$,
\begin{equation}\label{eq:touchard-power}
 B_{n+p^r}\equiv B_{n+1}+rB_n\pmod p.
\end{equation}
\end{theorem}
\begin{proof}
Work with remainder polynomials modulo $Q(X)=X^p-X-1$ over
$\mathbb F_p$. Each remainder has degree less than $p$ and is a finite list
of residues. The binomial theorem in characteristic $p$ says
$(U+V)^p=U^p+V^p$. Induction on $r$ therefore gives
\[
 X^{p^r}=X+r\quad\text{in }\mathbb F_p[X]/(Q):
\]
the case $r=0$ is immediate; raising to the $p$th power replaces $X+r$ by
$X^p+r=X+(r+1)$, using the elementary residue identity $r^p=r$.
Thus $X^{p^r}-X-r$ is a multiple of $Q$. Every shift of
\eqref{eq:touchard} says $(Q(E)B)_n=0$ modulo $p$; multiplying by a polynomial
and taking the corresponding finite linear combination of these shifts
preserves zero. This gives \eqref{eq:touchard-power}.
\end{proof}

\begin{theorem}[An explicit period modulo a prime]\label{thm:prime-period}
For every prime $p$, let
\[
  N_p=1+p+\cdots+p^{p-1}=\frac{p^p-1}{p-1}.
\]
Then, for all $n\geq0$,
\begin{equation}\label{eq:prime-period}
 B_{n+N_p}\equiv B_n\pmod p.
\end{equation}
The least positive period exists by bounded search and divides $N_p$.
\end{theorem}
\begin{proof}
In the same finite quotient ring,
\[
 X^{N_p}=\prod_{r=0}^{p-1}X^{p^r}
       =\prod_{r=0}^{p-1}(X+r)=X^p-X=1.
\]
The middle product equals $X^p-X$ by
\eqref{eq:finite-field-product}, since $r\mapsto-r$ permutes the residues.
Hence $Q$ divides $X^{N_p}-1$, and the shift calculation just used proves
\eqref{eq:prime-period}. No field-extension existence theorem or
irreducibility theorem is required: the ring consists explicitly of $p^p$
remainder lists and need not first be shown to be a field.

A candidate period $t$ can be tested by the $p$ congruences
$B_{t+i}\equiv B_i$ for $0\leq i<p$. The order-$p$ recurrence
\eqref{eq:touchard} propagates those congruences to every subsequent index.
Search $1\leq t\leq N_p$ for the least passing value $T$; the preceding
paragraph supplies a passing value. Divide $N_p=qT+s$, $0\leq s<T$.
Repeated use of period $T$ and period $N_p$ gives
$B_{n+s}\equiv B_n$ for every $n$. If $s>0$, it passes the same finite test,
contrary to minimality. Therefore $s=0$.
\end{proof}

For $p=2$, $N_p=3$ and the initial residues are $1,1,0$. Thus
\begin{equation}\label{eq:bell-parity}
 B_n\text{ is even}\quad\Longleftrightarrow\quad n\equiv2\pmod3.
\end{equation}
For $p=3$, the analogous finite test shows that the least period is $13$:
$13$ is prime, a period divides it, and period one fails since
$B_1\not\equiv B_2\pmod3$. The statement that $N_p$ is the least period for
every prime is not proved here; the proof establishes that it is a period
and that the least period divides it.

\subsection{Periodicity for every modulus}

\begin{theorem}[A primitive recursive period for every modulus]
\label{thm:all-moduli-period}
For every integer $m\geq1$, the Bell sequence is purely periodic modulo $m$.
There is a period $T$ with $1\leq T\leq m^m$, found by a bounded computation.
In particular the following uniform period is valid:
\begin{equation}\label{eq:all-moduli-period}
 B_{n+(m^m)!}\equiv B_n\pmod m\qquad(n\geq0).
\end{equation}
\end{theorem}
\begin{proof}
For every polynomial $f\in\mathbb Z[X]$, the binomial theorem gives
$f(X+m)\equiv f(X)\pmod m$, coefficientwise. Hence the finite shift identity
gives
\[
 \mathcal L\bigl(((X)_m-1)f(X)\bigr)
   =\mathcal L(f(X+m)-f(X))\equiv0\pmod m.
\]
Here $\mathcal L$ takes integer-coefficient polynomials to integers, because
its values on monomials are the integer Bell numbers. Write
\[
 Q_m(X)=(X)_m-1=X^m+c_{m-1}X^{m-1}+\cdots+c_1X-1.
\]
Taking $f=X^n$ yields the order-$m$ recurrence
\begin{equation}\label{eq:composite-recurrence}
 B_{n+m}\equiv B_n-\sum_{j=1}^{m-1}c_jB_{n+j}\pmod m.
\end{equation}
Its state vector is
$v_n=(B_n,\ldots,B_{n+m-1})\bmod m$, an element of a list of $m^m$ states.
The transition shifts the entries left and appends the right side of
\eqref{eq:composite-recurrence}. It is a bijection: from the new state,
recover the discarded entry by adding
$\sum_{j=1}^{m-1}c_jB_{n+j}$ to its last entry.

Among $v_0,\ldots,v_{m^m}$ two states agree, by the finite pigeonhole
principle. That principle here is proved by scanning the codes
$0,\ldots,m^m-1$: if no code repeats, the number of entries cannot exceed the
number of codes. If $v_a=v_b$ with $0\leq a<b\leq m^m$, apply the explicit
inverse transition $a$ times to get $v_0=v_{b-a}$. Therefore
$T=b-a$ is a period from the start; the forward recurrence proves
$v_{n+T}=v_n$ for every $n$. Such a $T$ is found by scanning
$1,\ldots,m^m$ for a return to $v_0$. Since every positive integer at most
$m^m$ divides $(m^m)!$, repeating the period proves
\eqref{eq:all-moduli-period}. The case $m=1$ is included, with its unique
residue state.
\end{proof}

The general bound is deliberately generous. Its purpose is to give a
uniform primitive recursive construction, while
Theorem~\ref{thm:prime-period} provides the much sharper bound at prime
moduli. The accompanying exact-arithmetic program certifies the source's
least periods for moduli $1,\ldots,12$:
\[
 1,3,13,12,781,39,137257,24,39,2343,28531167061,156.
\]
For every listed modulus other than $11$ it checks the complete orbit up to
its first return. For $11$ it uses the shorter polynomial certificate
\[
 N_{11}=28531167061=15797\cdot1806113,
 \qquad B_{15797}\equiv6,\quad B_{1806113}\equiv9\pmod{11}.
\]
Bounded trial division verifies primality of the two factors. Binary powering
with the reduction rule $X^{11}=X+1$ verifies the residues; the output includes
both remainder coefficient lists. Since $B_0\equiv1$, neither maximal proper
divisor is a period. The least period divides $N_{11}$, so it must equal
$N_{11}$. These are finite computation certificates with a specified checker,
not claims that the accompanying Python program is an \ipra\ proof checker.
Checks at selected moduli do not prove a universal least-period conjecture.

\paragraph{Arithmetic audit.}
Primality is the bounded test that no $2\leq d<p$ divides $p$. The proof of
Touchard's congruence uses at most $p$ iterates per finite partition code.
The quotient-ring proof uses lists of length $p$ with residue coefficients;
raising to $p^r$ can be replaced by $r$ successive Frobenius steps.
The general period proof scans at most $m^m+1$ state codes, each with $m$
entries. Existence of a period therefore comes with a primitive recursive
witness. The propagation arguments are quantifier-free inductions on the
decidable congruences or on their bounded conjunctions.

\section{Growth with explicit arithmetic error bounds}
\label{sec:growth}

The aim of this section is stronger than a statement that the usual
asymptotic argument is effective.  We give inequalities with explicit
errors, and, for the principal asymptotic equivalences, an explicit
primitive recursive modulus.  The logarithms, exponentials, square roots,
integrals, and the function $W$ in the displayed formulas denote the
represented quantities described in the foundations section.  A
non-strict inequality between them means the corresponding family of
rational inequalities with arbitrarily small prescribed slack.  No
quantification over an unrestricted space of real functions is needed.

Throughout, for $n\geq3$, put
\[
 r=W(n),\qquad a=e^r=\frac nr,\qquad
 b=n(r+1),\qquad L=\log n.
\]
Thus $r>1$, $r+\log r=L$, and
\begin{equation}\label{eq:growth-r-bracket}
 L-\log L\leq r\leq L.
\end{equation}
The lower inequality follows from $\log r\leq\log L$.  The function
$W(n)$ can be computed without deciding equality between real numbers.
For example, invert $x\mapsto xe^x$ on the rational interval $[0,n]$.
If a current interval has width $w$ and midpoint $m$, approximate
$me^m-n$ to error at most $w/8$.  An approximation greater than $w/4$
or less than $-w/4$ certifies which half contains the root.  Otherwise
the root is within $3w/8$ of $m$, since the derivative of $xe^x$ is at
least one for $x\geq0$.  Intersect with that smaller interval.  Each
step decreases the width by a factor at most $3/4$, so a predetermined
primitive recursive number of steps suffices for any requested error.
Existence here is the compatibility of these explicit brackets; it is
not an application of an unrestricted intermediate value principle.

\subsection{A logarithmic envelope requiring no saddle-point argument}

We first record the elementary factorial bounds
\begin{equation}\label{eq:elementary-factorial-log}
 m\log m-m+1\leq\log(m!)
 \leq m\log m-m+1+\log m\qquad(m\geq1).
\end{equation}
They follow by bounding the increasing function $\log x$ on each of
the intervals $[j,j+1]$ and adding the resulting finite inequalities.
The integral $\int_1^m\log x\,dx=m\log m-m+1$ is obtained from its
explicit antiderivative.  Equivalently, the same proof uses rational
Riemann sums and their explicit error bounds.

\begin{theorem}[An explicit logarithmic envelope]\label{thm:log-envelope}
For every integer $n\geq3$, let
\[
 H(n)=n\left(r-1+\frac1r\right),\qquad r=W(n).
\]
Then
\begin{equation}\label{eq:log-envelope}
 H(n)-r-4\leq\log B_n\leq H(n)+\log n.
\end{equation}
In particular,
\begin{equation}\label{eq:log-envelope-symmetric}
 \left|\frac{\log B_n}{n}-
 \left(r-1+\frac1r\right)\right|
 \leq\frac{\log n+4}{n}.
\end{equation}
\end{theorem}

\begin{proof}
All coefficients of the exponential generating function are
nonnegative.  Consequently, for every positive $x$,
\[
 B_n\leq \frac{n!}{x^n}\exp(e^x-1).
\]
Set $x=r$ and use the upper bound in
\eqref{eq:elementary-factorial-log}.  Because
$\log(n/r)=r$ and $a=n/r$, this gives
\[
 \log B_n\leq n\log n-n+1+\log n+a-1-n\log r
 =H(n)+\log n.
\]

For the lower bound take one positive term of Dobi\'nski's formula.
There is a small constructive point in selecting that term: we do not
take the floor of an arbitrary represented real.  Compute a rational
lower approximation $u$ satisfying $a-1/4\leq u\leq a$, and take the
ordinary rational floor $k=\lfloor u\rfloor$.  Then
\[
 a-\frac54<k\leq a,\qquad k\geq\frac a2,\qquad k\geq1.
\]
Define
\[
 g(x)=n\log x-x\log x+x.
\]
We have $g'(a)=0$, $g(a)=H(n)$, and, for $k\leq x\leq a$,
\[
 |g''(x)|=\frac n{x^2}+\frac1x
 \leq\frac{6r}{a}.
\]
Taylor's formula with its integral remainder therefore gives
\[
 0\leq g(a)-g(k)
 \leq\frac{3r}{a}\left(\frac54\right)^2
 =\frac{75r}{16e^r}<2.
\]
Here $re^{-r}\leq e^{-1}$ for $r\geq1$, and $e>5/2$ suffices for
the last numerical bound.  Dobi\'nski and
\eqref{eq:elementary-factorial-log} now give
\[
 \log B_n\geq n\log k-\log(k!)-1
 \geq g(k)-\log k-2\geq H(n)-r-4.
\]
This proves both assertions.
\end{proof}

\subsection{The de Bruijn logarithmic expansion}

\begin{theorem}[The displayed de Bruijn expansion, with constants]
\label{thm:debruijn}
Let $L=\log n\geq16$ and $\ell=\log L$.  Then
\begin{align}
 \left|\frac{\log B_n}{n}
 -\left(L-\ell-1+\frac{\ell+1}{L}
             +\frac{\ell^2}{2L^2}\right)\right|
 &\leq\frac{6\ell^3}{L^3}+\frac{L+4}{n}
 \label{eq:debruijn-sharp-error}\\
 &\leq\frac{7\ell}{L^2}.
 \label{eq:debruijn-displayed-error}
\end{align}
Thus, with the usual notation interpreted by these bounds,
\[
 \frac{\log B_n}{n}
 =\log n-\log\log n-1+
 \frac{\log\log n+1}{\log n}
 +\frac12\left(\frac{\log\log n}{\log n}\right)^2
 +O\!\left(\frac{\log\log n}{(\log n)^2}\right).
\]
The integer threshold $n\geq2^{24}$ is sufficient.
\end{theorem}

\begin{proof}
We give the error calculation, so the order symbol conceals no
ineffective threshold.  Set
\[
 r_0=L-\ell+\frac\ell L+
             \frac{\ell^2/2-\ell}{L^2},\qquad
 u=\frac{r_0-L}{L}.
\]
For $L\geq16$ one has $\ell\geq2$ and $0\leq-u\leq\ell/L\leq1/4$.
The convergent geometric-series integral for $\log(1+u)$ gives
\[
 \left|\log(1+u)-u+\frac{u^2}{2}\right|
 \leq\frac{|u|^3}{3(1-|u|)}
 \leq\frac{\ell^3}{2L^3}.
\]
Also
\[
 0\leq u+\frac\ell L\leq\frac{2\ell}{L^2},\qquad
 \left|\frac{u^2}{2}-\frac{\ell^2}{2L^2}\right|
 \leq\frac{2\ell^2}{L^3}.
\]
Substitution in $r_0+\log r_0-L$ cancels the terms through order
$L^{-2}$ and yields
\[
 |r_0+\log r_0-L|
 \leq\frac{(5/2)\ell^2+(1/2)\ell^3}{L^3}
 \leq\frac{3\ell^3}{L^3}.
\]
The derivative of $x+\log x$ is at least one on $x>0$, so
\begin{equation}\label{eq:W-expansion-error}
 \left|r-\left(L-\ell+\frac\ell L+
             \frac{\ell^2/2-\ell}{L^2}\right)\right|
 \leq\frac{3\ell^3}{L^3}.
\end{equation}

For completeness, the reciprocal term can be controlled without
substituting an unquantified expansion.  Put $v=(L-r)/L$.
By \eqref{eq:growth-r-bracket}, $0\leq v\leq\ell/L\leq1/4$, and
\[
 0\leq r-(L-\ell)=-\log(1-v)
 \leq\frac{4\ell}{3L}.
\]
The exact identity $(1-v)^{-1}=1+v+v^2/(1-v)$ therefore gives
\begin{equation}\label{eq:W-reciprocal-error}
 \left|\frac1r-\left(\frac1L+\frac\ell{L^2}\right)\right|
 \leq\frac{3\ell^2}{L^3}.
\end{equation}
Adding \eqref{eq:W-expansion-error} and
\eqref{eq:W-reciprocal-error} shows that the error in
$r-1+1/r$ is at most $6\ell^3/L^3$; the terms
$-\ell/L^2$ and $+\ell/L^2$ cancel.  Theorem
\ref{thm:log-envelope} proves
\eqref{eq:debruijn-sharp-error}.

Finally $\ell^2\leq L$ for $L\geq16$.  Also $e^L\geq L^5/120$,
so
\[
 \frac{L+4}{e^L}\leq\frac{150}{L^4}
 \leq\frac\ell{L^2}.
\]
This proves \eqref{eq:debruijn-displayed-error}.  The inequalities used
here are elementary monotonicity inequalities on $L\geq16$, with
their initial cases checked by the exponential series.
\end{proof}

\begin{corollary}[Principal logarithmic and root growth]
\label{cor:root-growth}
With explicit convergence rates supplied by
\eqref{eq:debruijn-sharp-error},
\[
 \frac{\log B_n}{n\log n}\longrightarrow1,
 \qquad
 \log\left(\frac{B_n^{1/n}\log n}{n}\right)
 \longrightarrow-1,
 \qquad
 \frac{B_n^{1/n}\log n}{n}\longrightarrow e^{-1}.
\]
In particular, the ordinary power series $\sum B_nx^n$ has zero
radius of convergence.  These logarithmic conclusions alone do not
imply a relative-error asymptotic formula for $B_n$; the next
subsections supply the additional estimate.
\end{corollary}

\begin{proof}
Divide the explicit expression in Theorem \ref{thm:debruijn} by $L$
for the first assertion.  Add $\ell-L$ for the second, and use the
exponential-series error bound for the third.  To see explicitly that
the rates are primitive recursive, one may bound $\log L\leq\sqrt L$
for $L\geq16$ and then choose $L$ larger than a fixed integer power of
the reciprocal requested error.  The threshold $n=2^{2M}$ ensures
$\log n\geq M$ because $\log2\geq1/2$.
For zero radius, the third assertion makes $B_n^{1/n}$ larger than
any prescribed positive rational number beyond an explicitly bounded
index.  Hence, for every nonzero rational $x$, the terms $B_nx^n$
fail the necessary condition of tending to zero.
\end{proof}

\subsection{The Gaussian constant and an effective Stirling estimate}

Two elementary analytic calculations used below are included to make
the normalization and the arithmetic content explicit.  First,
\begin{equation}\label{eq:gaussian-integral}
 \int_{-\infty}^{\infty}e^{-t^2/2}\,dt=\sqrt{2\pi},\qquad
 \int_T^\infty e^{-t^2/2}\,dt\leq\frac{e^{-T^2/2}}T
 \quad(T>0).
\end{equation}
For the tail bound, use $1\leq t/T$ and integrate
$t e^{-t^2/2}$.  For the constant, square the integral truncated to
$[-R,R]$.  The disk of radius $R$ is contained in that square, which
is contained in the disk of radius $\sqrt2R$.  Integration over
annuli, using their explicit areas, gives
\[
 2\pi(1-e^{-R^2/2})
 \leq\left(\int_{-R}^R e^{-t^2/2}\,dt\right)^2
 \leq2\pi.
\]
Both the annular computation and multiplication of the two integrals
are limits of finite Riemann sums with uniform elementary error
bounds.  Taking a predetermined $R$ large enough for the displayed
tail proves \eqref{eq:gaussian-integral}, without a completeness or
choice principle.

\begin{lemma}[Stirling's formula with an explicit remainder]
\label{lem:stirling-effective}
For every integer $m\geq1$ there is a represented number
$\vartheta_m$ with $0\leq\vartheta_m\leq1/(12m)$ such that
\begin{equation}\label{eq:stirling-effective}
 m!=\sqrt{2\pi m}\left(\frac me\right)^m e^{\vartheta_m}.
\end{equation}
\end{lemma}

\begin{proof}
Define
\[
 A_m=\log(m!)-(m+1/2)\log m+m.
\]
For $x=(2m+1)^{-1}$, the geometric series for
$\log((1+x)/(1-x))$ gives
\[
 A_m-A_{m+1}
 =(m+1/2)\log(1+1/m)-1
 =\sum_{j\geq1}\frac{x^{2j}}{2j+1}.
\]
This is positive and at most
\[
 \frac{x^2}{3(1-x^2)}=\frac1{12m(m+1)}.
\]
Consequently the sequence $A_m$ has a computable limit $c$, and
\[
 0\leq A_m-c\leq\frac1{12m}.
\]
This notation for the limit abbreviates the finite estimate
$0\leq A_m-A_M\leq1/(12m)-1/(12M)$ for $M\geq m$.

We identify $c$ by the elementary Wallis calculation.  Put
$I_j=\int_0^{\pi/2}\sin^j t\,dt$.  Integration by parts gives
$jI_j=(j-1)I_{j-2}$, with $I_0=\pi/2$ and $I_1=1$.
Since $0\leq\sin t\leq1$,
\[
 I_{2m+1}\leq I_{2m}\leq I_{2m-1}
 =\frac{2m+1}{2m}I_{2m+1}.
\]
The recurrence also gives the exact identities
\[
 I_{2m}=\frac\pi2\frac{\binom{2m}{m}}{4^m},
 \qquad I_{2m}I_{2m+1}=\frac\pi{2(2m+1)}.
\]
These inequalities, with error at most $1/(2m)$ in the ratio, imply
\[
 \sqrt m\,\frac{\binom{2m}{m}}{4^m}\longrightarrow\frac1{\sqrt\pi}.
\]
On the other hand, direct cancellation gives
\[
 A_{2m}-2A_m
 =\log\left(\sqrt m\,\frac{\binom{2m}{m}}{4^m}\right)
 -\frac12\log2.
\]
The left side tends to $-c$ with error at most $5/(24m)$.
The explicit Wallis inequalities identify the right limit as
$-\tfrac12\log(2\pi)$.  Thus $c=\tfrac12\log(2\pi)$, which proves
\eqref{eq:stirling-effective}.  Every limiting step in this proof has
the displayed reciprocal-integer error bound.
\end{proof}

\subsection{A quantitative saddle-point theorem}

Fourier coefficient extraction from the exponential generating
function gives
\begin{equation}\label{eq:growth-fourier}
 \frac{B_n}{\mathcal A_n}
 =\sqrt{\frac b{2\pi}}\int_{-\pi}^{\pi}e^{\Phi(\theta)}\,d\theta,
 \quad
 \mathcal A_n=\frac{n!\,e^{a-1}}{r^n\sqrt{2\pi b}},
 \quad
 \Phi(\theta)=e^{r e^{i\theta}}-a-in\theta.
\end{equation}
This identity needs no general complex integration theorem.  Expand
the entire function in its explicitly convergent power series on
$|z|=r$ and integrate each finite truncation.  The integral of
$e^{i(j-n)\theta}$ is zero for $j\ne n$ and is $2\pi$ for $j=n$.
The exponential-series tail bounds justify the passage to the
represented integral.  Although the integrand is complex, the
integral is real because the values at $\theta$ and $-\theta$ are
complex conjugates.

\begin{theorem}[A four-term relative-error estimate]
\label{thm:quantitative-saddle}
Suppose $r\geq1$, $T\geq1$, and $T\leq\sqrt b/r$.  Put
\[
 \epsilon=T^3\sqrt{\frac rn}.
\]
Then
\begin{align}\label{eq:saddle-error}
 \left|\frac{B_n}{\mathcal A_n}-1\right|
 \leq{}&e^\epsilon-1+e^{-T^2/2}
       +\frac{24}{T}e^{-T^2/48}
       +3\sqrt b\,e^{-n/(6r^2)}.
\end{align}
\end{theorem}

\begin{proof}
Set $\delta=T/\sqrt b$.  The first two derivatives of the phase at
zero are $\Phi'(0)=0$ and $\Phi''(0)=-b$.  For real $\theta$,
\[
 \left|\frac{d^3}{d\theta^3}e^{r e^{i\theta}}\right|
 \leq e^r(r+3r^2+r^3)\leq5nr^2.
\]
Taylor's formula therefore gives, on $|\theta|\leq\delta$,
\[
 \Phi(\theta)=-\frac{b\theta^2}{2}+E(\theta),
 \qquad |E(\theta)|\leq\frac56nr^2\delta^3\leq\epsilon.
\]
The difference between the normalized integral over this central
interval and the corresponding Gaussian integral is at most
$e^\epsilon-1$.  Completing the Gaussian integral to the whole line
costs at most $e^{-T^2/2}$, by
\eqref{eq:gaussian-integral}.

There are two tail regions.  On $|\theta|\leq1/r$ we have
$|r\sin\theta|\leq1$, $|\sin\theta|\geq|\theta|/2$, and
$e^{r\cos\theta}\geq a/2$.  Since $\cos u\leq1-u^2/3$ for
$|u|\leq1$,
\begin{align}
 \Re\Phi(\theta)
 &=(e^{r\cos\theta}-a)
   +e^{r\cos\theta}(\cos(r\sin\theta)-1)\notag\\
 &\leq-\frac{ar^2\theta^2}{24}
 \leq-\frac{b\theta^2}{48}.
 \label{eq:growth-middle-tail}
\end{align}
Thus the two pieces $\delta\leq|\theta|\leq1/r$ contribute at most
\[
 \frac2{\sqrt{2\pi}}\int_T^\infty e^{-u^2/48}\,du
 \leq\frac{24}{T}e^{-T^2/48}.
\]

For $1/r\leq|\theta|\leq\pi$, use
$1-\cos\theta\geq2\theta^2/\pi^2\geq\theta^2/5$.
Writing $s=r(1-\cos\theta)$, we get $s\geq1/(5r)$, and
\begin{equation}\label{eq:growth-far-tail}
 \Re\Phi(\theta)
 \leq-a(1-e^{-s})
 \leq-\frac a{6r}=-\frac n{6r^2}.
\end{equation}
We used $1-e^{-s}\geq s/(1+s)$, which follows from $e^s\geq1+s$.
The normalized length of the full integration interval is
$\sqrt{2\pi b}<3\sqrt b$.  This proves the final term in
\eqref{eq:saddle-error} and completes the proof.
\end{proof}

\begin{theorem}[A primitive recursive modulus for the saddle formula]
\label{thm:saddle-modulus}
For integers $q\geq1$ and
\begin{equation}\label{eq:saddle-modulus}
 n\geq N(q):=2^{64q^2},
\end{equation}
one has
\begin{equation}\label{eq:saddle-relative-modulus}
 \left|\frac{B_n}{\mathcal A_n}-1\right|\leq\frac5{8q}.
\end{equation}
The threshold is deliberately generous; its role is to give a simple
fully specified arithmetic modulus.
\end{theorem}

\begin{proof}
Use $T=16q$.  At and above the threshold,
\[
 \frac n{\log n}\geq2^{30}q^8,
 \qquad n\geq12(\log n)^3,
 \qquad n\geq24q.
\]
For example, $n/\log n$ is increasing for $n>e$; at $N(q)$ its
value is at least $2^{64q^2-6}/q^2$.  The inequality
$q^{10}\leq2^{10q}$ and $64q^2-36\geq10q$ verify the first
bound.  The second follows similarly from
$q^6\leq2^{6q}$ and $64q^2-18-6q\geq40$.  These are elementary
integer inequalities.

Since $r\leq\log n$, the first bound gives
$\epsilon=T^3\sqrt{r/n}\leq1/(8q)$ and also verifies
$T\leq\sqrt b/r$.  Therefore $e^\epsilon-1\leq1/(4q)$.
The next two terms in \eqref{eq:saddle-error} are each at most
$1/(8q)$: for the third, it suffices to check
$\exp(16q^2/3)\geq12$, already implied by the first three terms of
the exponential series.  Finally $r+1\leq n$, so $\sqrt b\leq n$;
and $n\geq12(\log n)^3$ gives
\[
 3\sqrt b\,e^{-n/(6r^2)}
 \leq3n e^{-2\log n}=\frac3n\leq\frac1{8q}.
\]
Adding the four bounds proves the assertion.
\end{proof}

\begin{corollary}[The Lambert $W$ equivalents on the source pages]
\label{cor:lambert-asymptotic}
As $n$ grows, with primitive recursive convergence moduli,
\begin{align}
 B_n&\sim\frac{n!\,e^{e^{W(n)}-1}}
                  {W(n)^n\sqrt{2\pi n(1+W(n))}},
 \label{eq:precise-saddle-equivalent}\\
 B_n&\sim\frac{n!\,e^{e^{W(n)}-1}}
                  {W(n)^n\sqrt{2\pi W(n)^2e^{W(n)}}},
 \label{eq:odlyzko-equivalent}\\
 B_n&\sim\frac1{\sqrt n}
          \left(\frac n{W(n)}\right)^{n+1/2}
          \exp\left(\frac n{W(n)}-n-1\right).
 \label{eq:wikipedia-equivalent}
\end{align}
For each of the last two expressions, call it $\mathcal C_n$,
the concrete assertion
\begin{equation}\label{eq:lambert-relative-modulus}
 q\geq1,\ n\geq2^{64q^2}\quad\Longrightarrow\quad
 \left|\frac{B_n}{\mathcal C_n}-1\right|\leq\frac1q
\end{equation}
holds.
\end{corollary}

\begin{proof}
The first formula is Theorem \ref{thm:saddle-modulus}.
For the second, its denominator replaces $\sqrt{r+1}$ by
$\sqrt r$, so its ratio with the first approximation is
$\sqrt{(r+1)/r}$.  At the stated threshold,
$L\geq32q^2$, and \eqref{eq:growth-r-bracket} gives
$r\geq L/2\geq16q^2$.  In particular,
\[
 0\leq1-\sqrt{\frac r{r+1}}\leq\frac1{2r}.
\]
For the last formula, substitute Lemma
\ref{lem:stirling-effective} into
\eqref{eq:precise-saddle-equivalent}.  If $R=B_n/\mathcal A_n$,
the exact ratio to its proposed right side is
\[
 R\,e^{\vartheta_n}\sqrt{\frac r{r+1}}.
\]
Here $|R-1|\leq5/(8q)$, $0\leq\vartheta_n\leq1/(12n)$, and
$1-1/(2r)\leq\sqrt{r/(r+1)}\leq1$.  The upper bound follows from
$e^{\vartheta_n}\leq1+1/(6n)$; the lower bound follows by
multiplying the nonnegative lower bounds.  Both differ from one by
less than $1/q$, proving \eqref{eq:lambert-relative-modulus}.
\end{proof}

\subsection{Successive ratios}

The logarithmic convexity established earlier is enough to obtain
successive-ratio growth without differentiating an asymptotic error
term.

\begin{corollary}[A quantitative ratio estimate]\label{cor:bell-ratio}
For every integer $n\geq16$,
\begin{equation}\label{eq:bell-ratio-error}
 \left|\log\frac{B_{n+1}}{B_n}-W(n)\right|
 \leq\frac{4\log n+22}{\sqrt n}.
\end{equation}
Consequently
\[
 \frac{B_{n+1}}{B_n}\sim\frac n{W(n)}\sim\frac n{\log n}.
\]
\end{corollary}

\begin{proof}
Let $F(n)=\log B_n$ and $h=\lfloor\sqrt n\rfloor$; the latter is
the integer square root, a bounded-search primitive recursive
operation.  Logarithmic convexity says that the increments of $F$
are nondecreasing.  Hence
\[
 \frac{F(n)-F(n-h)}h\leq F(n+1)-F(n)
 \leq\frac{F(n+h)-F(n)}h.
\]
Extend the explicitly defined $H(x)=x(W(x)-1+1/W(x))$ to positive
represented arguments.  Direct implicit differentiation of
$W(x)e^{W(x)}=x$ gives
\[
 H'(x)=W(x),\qquad
 0<H''(x)=\frac{W(x)}{x(1+W(x))}\leq\frac1x.
\]
Thus either secant slope of $H$ differs from $W(n)$ by at most
$h/(n-h)\leq2/\sqrt n$.  By
Theorem \ref{thm:log-envelope}, each endpoint error $F(j)-H(j)$,
for $n-h\leq j\leq n+h$, has absolute value at most
$\log n+5$.  Since $h\geq\sqrt n/2$, the error in either secant is
at most $(4\log n+20)/\sqrt n$.  This proves
\eqref{eq:bell-ratio-error}.  Exponentiation and
\eqref{eq:growth-r-bracket} give the two equivalents.
All uses of derivatives can equivalently be read as the displayed
finite increment bounds with integral remainders.
\end{proof}

\subsection{Why the growth proofs remain primitive recursive}

The preceding proof uses a fixed finite collection of analytic
functions on explicit bounded intervals, not a general theory of
real analysis.  Exponentials and trigonometric functions are evaluated
by finite Taylor sums with their elementary tails; logarithms use
the geometric-series integral after rational range reduction;
positive square roots and $W$ use bounded rational approximation.
All argument-size and precision bounds are primitive recursive.

For a particularly transparent integration bound, the phase in
\eqref{eq:growth-fourier} satisfies
$\Re\Phi(\theta)\leq0$ and $|\Phi'(\theta)|\leq2n$.
Therefore the integrand $e^{\Phi(\theta)}$ has Lipschitz constant
$2n$ on $[-\pi,\pi]$.  A uniform Riemann sum with $M$ intervals
has error, after the normalization in
\eqref{eq:growth-fourier}, at most $32n^2/M$.  Thus
$M\geq2^{t+6}n^2$ suffices for error at most $2^{-t}$; approximating
the finitely many endpoints and integrand values consumes a further
explicit portion of the error budget.  Gaussian tails are bounded by
\eqref{eq:gaussian-integral}, and all other tails were bounded in
the proofs above.  Consequently every integral equality or
inequality used here expands into rational finite-sum assertions at
a prescribed precision.

The arbitrary precision parameter is kept free during the finite
proof.  Inductions establish bounds for finite sums, products,
Taylor polynomials, and rational approximation programs; their
invariants are decidable arithmetic inequalities.  After all
represented expressions have been replaced by interval endpoints
and denominators cleared, these are quantifier-free induction
instances in the language with primitive recursive function
symbols.  The moduli in \eqref{eq:saddle-modulus},
\eqref{eq:debruijn-sharp-error}, and
\eqref{eq:bell-ratio-error} prevent any hidden use of an unbounded
search for an asymptotic threshold.  This is the intended iPRA
reading of the human proof; compiling these finite certificates
into the repository's object-language derivation type is a
separate formalization task, specified in the formalization section.


\section{The universal numerical upper bound}\label{sec:universal-bound}

The constant $0.792$ on the source page is not merely an asymptotic claim.
This section proves the exact all-positive-integer assertion of
Berend--Tassa~\cite{berend-tassa}, using a simple analytic estimate followed
by a small, completely rational certificate.  The estimate used here is
independent of the detailed saddle-point expansion.

\begin{theorem}[The $0.792$ bound]\label{thm:0792}
For every integer $n\ge1$,
\begin{equation}\label{eq:0792}
 B_n<\left(\frac{99n}{125\log(n+1)}\right)^n
     =\left(\frac{0.792n}{\log(n+1)}\right)^n.
\end{equation}
\end{theorem}
\begin{proof}
Take $r=\tfrac12\log(n+1)$ in
\eqref{eq:coefficient-bound}.  The elementary factorial estimate
\eqref{eq:elementary-factorial-log} gives
\begin{equation}\label{eq:0792-envelope}
 \frac{B_n^{1/n}\log(n+1)}n
 \le2\exp\left(-1+\frac{\log n+\sqrt{n+1}}n\right).
\end{equation}
The function $\log x/x$ decreases for $x\ge3$, since its derivative is
$(1-\log x)/x^2<0$.  The function $\sqrt{x+1}/x$ also decreases for $x>0$,
as direct differentiation gives
$-(x+2)/(2x^2\sqrt{x+1})<0$.
For $n\ge512$, therefore,
\[
 \frac{\log n+\sqrt{n+1}}n
 \le\frac{\log512+\sqrt{513}}{512}
 <\frac{7+23}{512}<\frac1{16}.
\]
Here $23^2>513$ and $e^7>512$ are elementary finite inequalities; for
example $e>5/2$ gives $(5/2)^7>512$.  Finally
\[
 e^{15/16}>\sum_{j=0}^{4}\frac{(15/16)^j}{j!}>\frac{250}{99}.
\]
The last comparison is a rational inequality verified by clearing
denominators.  Hence $2e^{-15/16}<99/125$, proving
\eqref{eq:0792} for every $n\ge512$.

The remaining finite assertion is
\begin{equation}\label{eq:small-bound-certificate}
 \bigwedge_{n=1}^{511}
 B_n\left(\frac{u_n}{2^{40}}\right)^n
       <\left(\frac{99n}{125}\right)^n,
 \qquad \log(n+1)\le\frac{u_n}{2^{40}}.
\end{equation}
We specify the integers $u_n$ by an exact bounded algorithm and supply all
511 values in \texttt{code/bound\_certificate.csv}.
Write $n+1=2^k y$ with $1\le y<2$, and $z=(y-1)/(y+1)$.  Define
\[
 L_J(z)=2\sum_{j=0}^{J-1}\frac{z^{2j+1}}{2j+1},\qquad
 U_J(z)=L_J(z)+\frac{2z^{2J+1}}{(2J+1)(1-z^2)}.
\]
With $J=24$, put
\[
 u_n=\left\lceil2^{40}\bigl(kU_J(1/3)+U_J(z)\bigr)\right\rceil.
\]
Formula \eqref{eq:log-series} proves the second inequality in
\eqref{eq:small-bound-certificate}.  The first is checked using only integers:
\begin{equation}\label{eq:certificate-integer-test}
 B_n(125u_n)^n<(99n\,2^{40})^n.
\end{equation}
The Bell recurrence generates $B_1,\ldots,B_{511}$.  The included verifier
performs precisely these 511 comparisons, without floating point arithmetic,
and all are true.  This is a finite computation proof: the arithmetic input,
all endpoints, and the complete algorithm are specified.  It can be expanded
into ground defining equations of \ipra, followed by their finite conjunction.
The Python run is independent reproducibility evidence, not a claim that a
Python interpreter checks the object logic.  Combining this finite certificate
with the already-proved infinite range proves the theorem.
\end{proof}

The near-maximal constant among these small indices occurs at $n=4$:
$15^{1/4}\log5/4$ is close to $0.792$.  This explains why ordinary
floating-point rounding is unsuitable as a proof of all the small cases.
The certificate uses an upper enclosure for each logarithm and an exact
strict integer comparison.

\begin{corollary}[An explicit version of the eventual page bound]
For every $n\ge2^{16}$,
\begin{equation}\label{eq:eventual-six-tenths}
 B_n<\left(\frac{e^{-3/5}n}{\log(n+1)}\right)^n.
\end{equation}
Consequently, for every positive rational $\eta$, the same threshold gives
$B_n<(e^{-3/5+\eta}n/\log(n+1))^n$.
\end{corollary}
\begin{proof}
Use $r=\tfrac34\log(n+1)$ in the coefficient bound.  The same cancellation
as before yields
\[
 \frac{B_n^{1/n}\log(n+1)}n
 \le\frac43\exp\left(-1+
             \frac{\log n+(n+1)^{3/4}}n\right).
\]
For $n\ge2^{16}$, the decreasing logarithm quotient is at most
$16/2^{16}$ because $\log2<1$.  Also
\[
 \frac{(n+1)^{3/4}}n
 =n^{-1/4}(1+1/n)^{3/4}
 \le\frac1{16}(1+2^{-16}).
\]
Their sum is less than $1/10$, by a rational comparison.
Since $e^{3/10}>1+3/10+(3/10)^2/2>4/3$, we have
$\log(4/3)<3/10$.  The displayed bound is therefore strictly less than
$\exp(-1+1/10+3/10)=e^{-3/5}$.
\end{proof}

The source page also supplies a particular threshold using an inverse
function $d^{-1}$.  The last corollary proves the same eventual inequality
with a different, explicit integer threshold and without computing that
inverse.  The sharper root-growth limit of Corollary~\ref{cor:root-growth}
further yields the optimal asymptotic constant $e^{-1}$: for every rational
$\eta>0$ there is a primitive recursively specified threshold beyond which
$B_n<(e^{-1+\eta}n/\log(n+1))^n$.

\section{Two correction terms, uniformly in a logarithmic shift}
\label{sec:corrections}

The longer expansion in the growth-rate section of Wikipedia can also be
obtained with an arithmetic remainder.  We give its coefficients explicitly
and bound the remainder, rather than invoke a general asymptotic-expansion
theorem.  This is the form of the expansion associated with
Moser--Wyman~\cite{moser-wyman}.

Put $r=W(n)$, $a=e^r$, and
\[
 U_j(x)=\sum_{k=0}^jS(j,k)x^k,\qquad
 U_j=U_j(r),\qquad v=U_2=r^2+r,\qquad b=av.
\]
For example $U_3=r^3+3r^2+r$ and
$U_4=r^4+6r^3+7r^2+r$.  Thus every $U_j$ below is an explicitly defined
polynomial, computable using the Stirling recurrence.
Define the following rational functions of $r$:
\begin{align*}
 P_0&=\frac{U_4}{8v^2}-\frac{5U_3^2}{24v^3},&
 P_1&=-\frac{U_3}{2v^2},&
 P_2&=-\frac1{2v},\\
 Q_0&=-\frac{U_6}{48v^3}+
       \frac{7U_3U_5}{48v^4}+\frac{35U_4^2}{384v^4}
       -\frac{35U_3^2U_4}{64v^5}+\frac{385U_3^4}{1152v^6},\\
 Q_1&=\frac{U_5}{8v^3}-\frac{35U_3U_4}{48v^4}
                        +\frac{35U_3^3}{48v^5},\\
 Q_2&=-\frac{5U_4}{16v^3}+\frac{35U_3^2}{48v^4},&
 Q_3&=\frac{5U_3}{12v^3},&
 Q_4&=\frac1{8v^2}.
\end{align*}
Write $P(r,h)=P_0+hP_1+h^2P_2$ and
$Q(r,h)=Q_0+hQ_1+h^2Q_2+h^3Q_3+h^4Q_4$.

\begin{theorem}[A uniform two-correction expansion]\label{thm:two-corrections}
Let $H\ge0$ be an integer, $n\ge3$, and $h\in\mathbb Z$ satisfy
$n+h\ge0$ and $|h|\le Hr$.  If $a\ge2^{12}$, then
\begin{equation}\label{eq:two-corrections}
 B_{n+h}=\frac{(n+h)!\,e^{a-1}}{r^{n+h}\sqrt{2\pi av}}
    \left(1+\frac{P(r,h)}a+\frac{Q(r,h)}{a^2}
                        +\mathcal E_{n,h}\right),
 \qquad
 |\mathcal E_{n,h}|\le\frac{2^{400}(H+1)^6}{a^3}.
\end{equation}
In particular this gives the displayed two-correction expansion uniformly for
$h=O(\log n)$, because $\log n=r+\log r\le2r$ for $r\ge1$.
\end{theorem}

\begin{proof}
Circular coefficient extraction at the same radius $r=W(n)$ gives
\[
 \frac{B_{n+h}r^{n+h}\sqrt{2\pi av}}{(n+h)!e^{a-1}}
 =\sqrt{\frac{av}{2\pi}}\int_{-\pi}^{\pi}
       \exp\bigl(\Phi(\theta)-ih\theta\bigr)\,d\theta,
\]
where $\Phi$ is the phase in \eqref{eq:growth-fourier}.  Set
$\varepsilon=a^{-1/2}$, $t=\sqrt{av}\,\theta$, and
$T=a^{1/12}$.  We analyze $|t|\le T$ and the two tails separately.

\paragraph{Taylor expansion of the phase.}
Induction on $j$ using the chain rule gives
\[
 \frac{d^j}{d\theta^j}e^{r e^{i\theta}}
   =i^j e^{r e^{i\theta}}U_j(r e^{i\theta}),
 \qquad
 \left|\frac{d^j}{d\theta^j}e^{r e^{i\theta}}\right|
   \le aU_j(r).
\]
The polynomial recurrence is $U_{j+1}(x)=x(U_j(x)+U_j'(x))$.
Since $r\ge1$ and $v\ge r^2$,
\begin{equation}\label{eq:touchard-derivative-bound}
 U_j(r)/v^{j/2}\le B_j.
\end{equation}
The explicit values $B_3=5,B_4=15,B_5=52,B_6=203,B_7=877,B_8=4140$
satisfy $B_j<j!$ for $3\le j\le8$.
Taylor's estimate through order seven therefore gives
\[
 \Phi(\theta)-ih\theta=-t^2/2+
       \sum_{j=1}^5\varepsilon^j A_j(t)+R(t),
 \qquad |R(t)|\le\varepsilon^6|t|^8,
\]
where
\begin{align*}
 A_1&=-i\left(\frac{ht}{\sqrt v}+
                         \frac{U_3t^3}{6v^{3/2}}\right),&
 A_2&=\frac{U_4t^4}{24v^2},\\
 A_3&=\frac{iU_5t^5}{120v^{5/2}},&
 A_4&=-\frac{U_6t^6}{720v^3},&
 A_5&=-\frac{iU_7t^7}{5040v^{7/2}}.
\end{align*}
In particular $|A_1|\le H|t|+|t|^3$ and
$|A_j|\le|t|^{j+2}$ for $2\le j\le5$.

\paragraph{An explicit sixth-order remainder for the exponential.}
Fix $t$ and put $Z(x)=\sum_{j=1}^5x^jA_j(t)$ for $0\le x\le\varepsilon$.
Since $a\ge2^{12}$, we have
$\varepsilon T=a^{-5/12}\le1/16$.  The only real parts in $Z$ come from
$A_2$ and $A_4$, and $A_4\le0$, so
\[
 \Re Z(x)\le x^2t^4\le t^2/256.
\]
For $1\le k\le5$,
\[
 |Z^{(k)}(x)|\le D(t):=600(H+1)(1+|t|^7).
\]
Indeed, every derivative coefficient of a monomial of degree at most five
is at most $5!=120$, $x\le1$, and
$H|t|+|t|^3+\cdots+|t|^7\le5(H+1)(1+|t|^7)$.
The sixth derivative of $e^Z$ is $e^Z$ times the complete Bell polynomial
in $Z',\ldots,Z^{(6)}$.  This follows by six applications of the product
rule, or from \eqref{eq:block-multiplicity}; its coefficient sum is
$B_6=203$, and every monomial has at most six factors.  Since $D\ge1$,
\[
 |(e^Z)^{(6)}|\le203D^6e^{t^2/256}.
\]
Let $C_j(t)$ be the coefficient of $x^j$ in the degree-five Taylor polynomial
of $e^{Z(x)}$ at zero.  Taylor's integral remainder yields
\[
 \left|e^{Z(\varepsilon)}-\sum_{j=0}^5\varepsilon^j C_j(t)\right|
 \le\varepsilon^6\frac{203}{720}D^6e^{t^2/256}.
\]
The additional phase remainder contributes at most
$\varepsilon^6|t|^8\exp(t^2/128)$, since
$|R|\le t^2(\varepsilon|t|)^6\le t^2/256$ and
$|e^R-1|\le|R|e^{|R|}$.  After multiplication by $e^{-t^2/2}$, the sum of
these two error bounds is at most
\begin{equation}\label{eq:mw-integrable-remainder}
 \varepsilon^6 2^{66}(H+1)^6(1+|t|^{42})e^{-t^2/4}.
\end{equation}
For example, use $600^6<2^{60}$,
$(1+|t|^7)^6\le2^5(1+|t|^{42})$, and
$|t|^8\le1+|t|^{42}$.
The normalized full-line integral of
$(1+|t|^{42})e^{-t^2/4}$ is
$\sqrt2(1+2^{21}41!!)<2^{149}$.
Thus the central integrated error is less than
$2^{220}(H+1)^6a^{-3}$.

\paragraph{Coefficient evaluation.}
The odd $C_j$ are imaginary odd polynomials, so their integrals on the
symmetric central interval vanish exactly.  The even coefficients needed are
\[
 C_0=1,\quad C_2=A_2+\tfrac12A_1^2,
 \quad
 C_4=A_4+A_1A_3+\tfrac12A_2^2+\tfrac12A_1^2A_2+\tfrac1{24}A_1^4.
\]
The full-line Gaussian averages of $C_2,C_4$ are $P(r,h),Q(r,h)$,
respectively.  Substituting the $A_j$ and using the even moments
$1,3,15,105,945,10395$ proves exactly the displayed coefficients.
These moments satisfy $\mu_{2j}=(2j-1)\mu_{2j-2}$ by integration by parts,
with $\mu_0=1$; all odd moments vanish.
This is a finite polynomial calculation, including all terms in $h$.

Replacing the central Gaussian polynomial integrals by full-line integrals
also has a bounded error.  The same derivative bound gives
\[
 |C_0|+|C_2|+|C_4|
 \le2^{66}(H+1)^6(1+|t|^{42}).
\]
On $|t|\ge T$, multiply the integrand by the upper bound
$|t|^{36}/T^{36}$ for $1$.  Since $T^{36}=a^3$, the resulting error is at most
\[
 \frac{2^{66}(H+1)^6}{a^3}
  \frac1{\sqrt{2\pi}}
       \int_{\mathbb R}(|t|^{36}+|t|^{78})e^{-t^2/2}\,dt
 <\frac{2^{350}(H+1)^6}{a^3}.
\]
Here $77!!\le78^{39}<2^{273}$ and the smaller moment is harmless.

\paragraph{The noncentral part of the original integral.}
The factor $e^{-ih\theta}$ has absolute value one, so the tail bounds
\eqref{eq:growth-middle-tail} and \eqref{eq:growth-far-tail} are unchanged.
The middle tail is at most $(24/T)e^{-T^2/48}$.  Using
$e^x\ge x^{18}/18!$ shows that this is at most
$24\cdot18!\cdot48^{18}a^{-3}<2^{170}a^{-3}$.
For the far tail, the elementary inequality $r\le e^{r/2}=\sqrt a$
implies $\sqrt b\le2a$ and $a/(6r)\ge\sqrt a/6$.
Consequently
\[
 3\sqrt b\,e^{-a/(6r)}\le6a e^{-\sqrt a/6}
 \le 8!\,6^9 a^{-3}<2^{44}a^{-3}.
\]
The inequality $r\le e^{r/2}$ follows already from
$1+r/2+r^2/8\ge r$.  The required central cutoff lies inside
$|\theta|\le1/r$, because $T=a^{1/12}\le\sqrt{av}/r$.

Adding the central error, the Gaussian polynomial tail, and the two original
tails gives less than $2^{400}(H+1)^6a^{-3}$.
This proves \eqref{eq:two-corrections}.
\end{proof}

Every constant in this proof is intentionally generous.  The uniform
remainder is the main point: it verifies the two orders of correction and
their dependence on $h$ without hiding an effectiveness hypothesis in an
$O$-symbol.  For an explicit accuracy request $1/q$, set
$C_H=2^{400}(H+1)^6$ and take
\[
 n\ge\max\{2^{24},(C_H(q+1))^2\}.
\]
The inequality $r\le\sqrt a$ gives $n=ar\le a^{3/2}$ and hence
$a\ge\sqrt n$ when $a\ge1$.  Thus this threshold ensures both
$a\ge2^{12}$ and $C_H/a^3\le1/q$ for $q\ge1$.
Fractional powers and the integrations can be replaced by the root and
Riemann-sum algorithms of Section~\ref{sec:foundations}; all bounds remain
primitive recursive in $H$ and $q$.

\section{Lean implementation and the route to object-theory proofs}
\label{sec:formalization}

\subsection{What the supplied Lean files establish}

The accompanying sources were compiled with Lean~4.32.0, using only the
bundled \texttt{Std} library. The file \texttt{BellCore.lean} defines finite
sums, the two-variable Stirling recurrence, Bell numbers as row sums, and
factorials. It proves the Stirling boundary identities, positivity of Bell
numbers, the weighted recurrence
\[
 B_{n+1}=\sum_{k=0}^{n}(k+1)S(n,k),
\]
monotonicity, the upper bounds $B_{n+1}\le(n+1)B_n$ and $B_n\le n!$, and
the values through $B_{10}=115975$. The numerical identities use
definitional reduction. The file \texttt{Audit.lean} prints the trusted-axiom
dependencies of nine representative theorems. The computational identities
and the column-one identity use no axioms; the other printed proofs use
Lean's \texttt{propext} and \texttt{Quot.sound}. None uses
\texttt{Classical.choice} or \texttt{sorryAx}.

These facts describe the delivered formalization exactly. In particular,
the package does \emph{not} contain a Lean-checked proof of the saddle-point
theorem, nor a constructed \texttt{ConciseIPRA.Prf} derivation of the Bell
theory. Ordinary Lean theorems about natural numbers are metatheorems, even
when their proofs are constructive. Constructing derivations in a specified
first-order theory is a separate task. The mathematics of the preceding
sections supplies the human proof; this section specifies the remaining
formalization work and supplies executable definitions for its main targets.

The companion file \texttt{BellTargets.lean} defines total arithmetic
programs and propositions with names ending in \texttt{Target}. Such a definition records a
goal; it does not assert that goal. There are no added axioms or placeholder
proofs. This distinction also applies to the proposed mixed-radix Bell
program and the approximation programs below: their definitions compile,
while their mathematical correctness proofs remain to be translated into
Lean.

\subsection{The inspected ProveIt interfaces}

The repository was inspected at commit
\texttt{c0d0dc37f60aea10c0182ba4eece9eae8c7e13c6}.
The relevant namespace is \texttt{LeanProofs.ConciseIPRA}.
The source root is
\begin{quote}\small
\path{Logic/PrimitiveRecursiveArithmetic/ConciseIPRA/Lean/ConciseIPRA}.
\end{quote}
The audit JSON records the Git blob identifiers of every inspected module.
The upstream project itself was not rebuilt for this package.

The
\href{https://github.com/VladimirReshetnikov/ProveIt/blob/c0d0dc37f60aea10c0182ba4eece9eae8c7e13c6/Logic/PrimitiveRecursiveArithmetic/ConciseIPRA/Lean/ConciseIPRA/PR.lean}{\texttt{PR.lean} interface}
contains the following exact objects.
\begin{itemize}
\item \texttt{PR k}: primitive-recursive symbols of arity $k$, generated by
\texttt{zero}, \texttt{succ}, \texttt{proj}, \texttt{comp}, and
\texttt{prec}. In the last constructor the recursion variable is the first
argument, and the step receives the counter, the previous value, and the
parameters.
\item \texttt{PTm n}: terms with $n$ free variables; substitution is
\texttt{PTm.sub} and its generic interface is \texttt{Syn}.
\item \texttt{PRAAx}: the defining equations, the two successor axioms, and
induction for formulas satisfying \texttt{Fm.QF}.
\item \texttt{PRA.precZero}, \texttt{PRA.precSucc}, and
\texttt{PRA.qfInd}: reusable constructors for the corresponding
object-theory derivations.
\end{itemize}
The
\href{https://github.com/VladimirReshetnikov/ProveIt/blob/c0d0dc37f60aea10c0182ba4eece9eae8c7e13c6/Logic/PrimitiveRecursiveArithmetic/ConciseIPRA/Lean/ConciseIPRA/Deduction.lean}{\texttt{Deduction.lean} interface}
defines \texttt{Prf Ax} as intuitionistic natural deduction with equality,
including hypotheses, instantiated theory axioms, introduction and
elimination rules for the connectives and quantifiers, and simultaneous
equality substitution. Its derived operations include \texttt{Prf.subst},
\texttt{Prf.trans}, \texttt{Prf.eqCongr}, and \texttt{Prf.eqTrans}.

The
\href{https://github.com/VladimirReshetnikov/ProveIt/blob/c0d0dc37f60aea10c0182ba4eece9eae8c7e13c6/Logic/PrimitiveRecursiveArithmetic/ConciseIPRA/Lean/ConciseIPRA/BiInterpretation.lean}{\texttt{BiInterpretation.lean} interface}
provides \texttt{pra\_iff\_teq} and \texttt{pra\_theorem\_iff}; the
\href{https://github.com/VladimirReshetnikov/ProveIt/blob/c0d0dc37f60aea10c0182ba4eece9eae8c7e13c6/Logic/PrimitiveRecursiveArithmetic/ConciseIPRA/Lean/ConciseIPRA/NativeBiInterpretation.lean}{native interface}
provides \texttt{pra\_iff\_native}. These translate an already constructed
$\ipra$ derivation to the concise scalar or native calculus. They do not
turn a semantic Bell-number theorem into a derivation automatically. The
iterative implementation can also use \texttt{Tm.prc} and its proved
equations \texttt{Q.prc\_zero}, \texttt{Q.prc\_succ} from
\texttt{TeqRec.lean}.

\subsection{First-order codes and bounded invariants}

All finite objects may be encoded by natural numbers. For definiteness,
one may use
\[
 \langle a,b\rangle=2^a(2b+1)-1.
\]
The projections search only below the input code and are primitive
recursive. Encode a signed integer by a sign and a magnitude, a rational
by a signed numerator and a positive denominator, and a finite array by
its length together with a mixed-radix word. Addition, multiplication,
comparison, division with remainder, factorial, powers, bounded sums,
bounded products, and bounded searches then have fixed PR definitions.
Invalid inputs receive specified default values, so every function is total.
Theorems about valid codes carry validity hypotheses.

An especially direct first-order Bell program fixes the input $N$ and sets
\[
 A=(N+2)^{N+2}+1,\qquad w=N+1,\qquad
 d_k(c)=\left\lfloor c/A^k\right\rfloor\bmod A.
\]
Start with $c_0=1$, and perform $N$ steps
\[
 c_{j+1}=\sum_{k=0}^{N}
 \begin{cases}
 0,&k=0,\\
 k\,d_k(c_j)+d_{k-1}(c_j),&k>0
 \end{cases} A^k.
\]
Finally return $\sum_{k=0}^{N}d_k(c_N)$. These are the definitions
\texttt{digit}, \texttt{packedStep}, \texttt{packedRows}, and
\texttt{packedBell} in the delivered source.

The invariant at stage $j\le N$ says that the digits are $S(j,k)$ for
$0\le k\le N$, and that every digit is below $A$. There is no carry because
\[
 S(j,k)\le B_j\le j!\le (N+1)^{N+1}<A.
\]
The invariant is a bounded conjunction of elementary comparisons and
equalities. Replace that conjunction by its primitive-recursive Boolean
test and apply induction to the single equation saying that the test is
$1$. When the final unbounded symbol $S(n,k)$ is defined from a table whose
size is determined by $n+k$, first prove that enlarging the table width or
the base does not change earlier entries. This is a second bounded
simultaneous invariant. It yields the usual Stirling recursion with no
appeal to higher-type recursion.

This construction matters because the elegant Lean definition
\texttt{stirling : Nat $\to$ Nat $\to$ Nat} recurses over a row viewed as a
function. Its ordinary Lean proof of vanishing above the diagonal also
generalizes over the second argument during induction. Neither feature
should be silently identified with quantifier-free induction in the
object theory. The packed-table invariant is the explicit replacement.

For the combinatorial interpretation, a restricted-growth word of length
$n$ is encoded in base $n+1$ below $(n+1)^n$. Its validity test checks the
initial label and the bounded inequalities against preceding maximum
labels. The number of blocks is the maximum label plus one, with the empty
word treated separately. Counting valid words, and counting those with
exactly $k$ labels, are therefore bounded sums. Every bijection used in the
finite proofs becomes an explicitly bounded encoder and decoder together
with their inverse equations. Finite groups in the congruence argument are
handled in exactly this fashion: their actions are permutations of bounded
code sets, and orbit calculations are bounded arithmetic.

There is no uniform primitive-recursive evaluator for \emph{all} PR
programs assumed here. Each displayed algorithm is compiled to one fixed
PR definition; subsequent formulas contain that symbol. Generic lemmas
about a function $f$ in Lean are metatheoretic schemas, instantiated with
these fixed programs before constructing an object-theory proof.

\subsection{Rational approximation programs and their error estimates}
\label{subsec:lean-approximation}

The record \texttt{QCode} consists of an integer numerator and a natural
denominator. Its order test is cross multiplication. The operations do not
require normalization, and positivity of a denominator is a separate
arithmetic property. Write $\varepsilon_t=2^{-t}$.

\paragraph{Exponential.}
For a valid rational code $x=u/v$, put $m=|u|+1$ and
\[
 K=8(m+1)^2+2t+16,\qquad
 E_t(x)=\sum_{j=0}^{K-1}\frac{x^j}{j!}.
\]
This is \texttt{expApprox}. It satisfies
\[
 |E_t(x)-e^x|\le 2^{-t}.
\]
Here and below the displayed real notation abbreviates the corresponding
finite Cauchy-tail inequalities. To prove the exponential bound, note that
$|x|\le m$ and, for $k\ge K$, the last half of the factors in $k!$ are at
least $\lfloor k/2\rfloor+1$. For $k=2h$ this gives
$m^k/k!\le(m^2/(h+1))^h\le 2^{-k}$; for $k=2h+1$ it gives the bound
$2^{1-k}$. Thus every finite tail starting at $K$ is bounded by
$\sum_{k\ge K}2^{1-k}=2^{2-K}\le2^{-t}$.
For nonnegative $x$, all summands are nonnegative, so $E_t(x)\le e^x$.
This one-sided property is used in the bounded search for $W$.

\paragraph{Logarithm.}
For $x=u/v>0$, with positive integers $u,v$, put
\[
 z=\frac{u-v}{u+v},\qquad
 M=(u+v)^2 2^{t+4},\qquad
 L_t(x)=2\sum_{j=0}^{M-1}\frac{z^{2j+1}}{2j+1}.
\]
This is \texttt{logApprox}. Its finite-tail bound is
\[
 \left|2\sum_{j=M}^{J}\frac{z^{2j+1}}{2j+1}\right|
 \le\frac{2}{(2M+1)(1-z^2)}
 =\frac{(u+v)^2}{2uv(2M+1)}
 \le2^{-t}.
\]
The logarithm identities used elsewhere follow by the ordinary finite
power-series manipulations with this tail bound, or by rational
integration of $1/x$. Invalid and nonpositive inputs are sent to zero in
the total program; no theorem interprets those outputs as logarithms.

\paragraph{Lambert $W$.}
Let $n\ge1$, $\delta=2^{-(t+2)}$, and test the dyadic grid
$x_k=k\delta$ for $0\le k\le(n+1)/\delta$. Use
\[
 s=t+n+8,
 \qquad x_kE_s(x_k)\ge n
\]
as the success test. Choose its first successful index by the explicit
bounded search \texttt{boundedMin}. The last test succeeds because
$E_s(n+1)\ge1$. Write $r=W(n)$ and $r_t$ for the selected grid point.
Since $E_s(x)\le e^x$ at nonnegative arguments, $r_t\ge r$.
At the preceding grid point the test fails, while its error in $xe^x$ is
at most
\[
 (n+1)2^{-s}\le2^{-t-8}<\delta/4.
\]
The function $F(x)=xe^x$ has slope at least $1$ on $[0,\infty)$; the
inequality $F(y)-F(x)\ge y-x$ also follows directly from the positive
power series for $F$. Consequently
\[
 0\le r_t-r<\delta+\delta/4<2^{-t}.
\]
All tests here compare finite rational numbers. The algorithm never asks
whether two represented real numbers are equal, or on which side of a
rational a possibly equal real lies. The function \texttt{wApprox} is
therefore a total PR program with the indicated modulus.

\paragraph{Square root and $\pi$.}
The source uses
\[
 R_t(x)=2^{-t}\left\lfloor
 \sqrt{\lfloor x2^{2t}\rfloor}\right\rfloor
\]
for nonnegative $x$. The integer square root may be implemented by bounded
search; its defining square inequalities give
$0\le\sqrt{x}-R_t(x)<2^{-t}$ directly. Lean's \texttt{Nat.sqrt} is used
for this executable presentation, with a bounded-search implementation
available for its PR compilation. The rational approximation to $\pi$ is
the first $2^{t+4}$ terms of
$4\sum_{k\ge0}(-1)^k/(2k+1)$. Pairing adjacent terms proves the alternating
tail estimate $4/(2^{t+5}+1)<2^{-t}$ by finite arithmetic. The usual
identification with $4\arctan(1)$ fixes the normalization of $\pi$.

These choices are deliberately generous. They make the PR bounds explicit
and easy to inspect; they are not intended as efficient numerical
implementations.

\subsection{Exact arithmetic targets for the asymptotic assertions}

The definition \texttt{LogLambertTarget} is the statement that, for all
$n\ge3$ and all $t$, the rational approximation to
\[
 \left|\frac{\log B_n}{n}-
 \left(r-1+\frac1r\right)\right|
\]
is bounded by $(L_t(n)+4)/n+8\varepsilon_t$, using $r_t$ in place of $r$.
This is a first-order arithmetic statement: every approximate quantity is
a fixed PR term and every comparison clears positive denominators. The
extra $8\varepsilon_t$ accounts for the approximation errors. Indeed,
$r_t\ge r\ge1$ gives
$|1/r_t-1/r|\le\varepsilon_t$, and the logarithmic estimates contribute
at most another $2\varepsilon_t/n$.

For the leading equivalent, use the algebraically simpler expression
\[
 C(n)=\left(\frac nr\right)^n
       \frac{\exp(n/r-n-1)}{\sqrt r}.
\]
It is exactly the Wikipedia/Lov\'asz expression, after cancelling the two
square roots involving $n$. The source \texttt{growthApprox n t} evaluates
this expression at internal precision
$s=t+8(n+1)^2+16$.
Here is an explicit error check for that overprecision. For $n\ge3$,
$r\ge1$, $r_s\ge r$, and $a_s=n/r_s\le a=n/r\le n$ satisfy
$|a_s-a|\le n\varepsilon_s$. Therefore
\[
 |a_s^n-a^n|\le n^{n+1}\varepsilon_s.
\]
Both exponential arguments are negative; the Lipschitz constant of the
exponential there is at most $1$. Thus its approximated value differs
from $\exp(a-n-1)$ by at most $(n+1)\varepsilon_s$, and has magnitude at
most $2$. The approximated square root is at least $1$ and differs from
$\sqrt r$ by at most $3\varepsilon_s/2$. Splitting the difference of the
three factors gives
\[
 |C_t(n)-C(n)|
 \le4(n+1)n^n\varepsilon_s
 \le2^{2n^2}\varepsilon_s
 <2^{-t}.
\]
The elementary inequalities used here all hold for $n\ge3$; in particular
$4(n+1)\le2^{n^2}$ and $n^n\le2^{n^2}$.

To control the final quotient, observe that
\[
 \log C(n)=n(r-1+1/r)-1-\tfrac12\log r
 \ge\tfrac12(n-1)\ge1.
\]
We used $r-1+1/r\ge1$, $\log r\le r-1$, and $r\le n$.
Thus $C(n)>2$. Re-evaluate $C$ at precision
$t+4(n+1)^2+8$; its rational approximation is at least $1$, and
$B_n\le n!\le n^n\le2^{n^2}$ gives
\[
 \left|\frac{B_n}{C_{t+4(n+1)^2+8}(n)}-
              \frac{B_n}{C(n)}\right|<2^{-t}.
\]
Accordingly \texttt{WikipediaGrowthTarget} states the rational inequality
with right side $1/q+4\varepsilon_t$, for
\[
 q\ge1,\qquad n\ge2^{64q^2}.
\]
The saddle-point theorem in the human proof supplies $1/q$; the preceding
finite error calculation supplies the rational slack. This is the precise
arithmetic replacement for the asymptotic equivalence. No unrestricted
quantification over real numbers is being added to $\ipra$.

For the two-correction expansion, the additional definitions are the finite
polynomials $U_j$ for $j\le8$, the rational expressions $P,Q$ displayed in
Section~\ref{sec:corrections}, and the normalized expression
\[
 D(n,h)=\frac{B_{n+h}r^{n+h}\sqrt{2\pi av}}{(n+h)!e^{a-1}}.
\]
The target is the finite-precision form of
$|D(n,h)-1-P(r,h)/a-Q(r,h)/a^2|\le2^{400}(H+1)^6/a^3$,
with the stated hypotheses. A signed integer code represents $h$.
All its extra algebra is finite polynomial multiplication and Gaussian
moments through degree $78$. The proof requires the eighth-order phase
Taylor bound and the sixth derivative of the exponential of a degree-five
polynomial, exactly as displayed in that section. These additional targets
are specified in the paper; unlike the two leading-growth targets, they
are not yet packaged as named definitions in \texttt{BellTargets.lean}.

\subsection{A concrete formalization order}

The outstanding work has the following dependency order.
\begin{enumerate}
\item Implement the fixed PR codes for integer/rational arithmetic,
finite arrays and bounded folds. Prove their defining equations and
positive-denominator closure using \texttt{PRA.precZero},
\texttt{PRA.precSucc}, equality congruence and \texttt{PRA.qfInd}.
\item Prove the packed-table invariant and
\texttt{PackedAgreementTarget}. Compile each finite-sum lemma as a schema
for fixed PR summands. Replace any ambient Lean induction over a
universally quantified family by its bounded Boolean invariant.
\item Formalize the bounded restricted-growth encodings and the explicit
bijections in the paper. Derive the Stirling and Bell recurrences, the
finite summation identities, and the count interpretations.
\item Formalize factorial cancellation and the finite moment identities.
For the determinant, use an explicit permutation-sum determinant and the
triangular change of polynomial basis. For congruences, use remainder
arithmetic, bounded orbit decompositions and the finite shift-polynomial
calculation. The propositions \texttt{TouchardTarget},
\texttt{PrimePeriodTarget}, \texttt{LogConvexTarget} and
\texttt{NormalizedLogConcaveTarget} give principal endpoints.
\item Prove the rational tail estimates above. Represent exponential,
logarithm, trigonometric functions and the particular square roots and
Lambert values by these fixed PR approximation programs. Prove their
algebraic and order laws at finite precision before using real notation
as an abbreviation.
\item For each integral in the analytic proof, write the finite rational
mesh sum, its explicit mesh error, and its truncation error. Prove finite
coefficient orthogonality, the local Taylor bounds and the two tail
bounds. Combine the estimates with the explicit thresholds from the
growth section. This yields the arithmetic targets just stated.
\item Only after constructing these \texttt{Prf PRAAx} derivations, apply
\texttt{pra\_iff\_teq} or \texttt{pra\_iff\_native} to transport them to
the selected concise calculus.
\end{enumerate}

The essential remaining workload is explicit: it is the construction of
the arithmetic derivations, especially the finite-precision integral
estimates. The supplied definitions, rates, and human proofs specify what
those derivations must prove. Neither an unrestricted completeness theorem
nor an unproved transfer from classical real analysis is used to declare
the task complete.

\clearpage
\section{Exact numerical certificates}\label{sec:numeric}

The program \texttt{code/verify\_numeric.py} computes the Bell recurrence
\eqref{eq:bell-binomial} through index $1000$, using integer arithmetic only.
Its output is \texttt{code/numeric\_validation.json}. The first values are:
\begin{center}
\begin{tabular}{r r@{\qquad}r r}
\hline
$n$&$B_n$&$n$&$B_n$\\
\hline
0&1&8&4140\\
1&1&9&21147\\
2&2&10&115975\\
3&5&11&678570\\
4&15&12&4213597\\
5&52&13&27644437\\
6&203&14&190899322\\
7&877&15&1382958545\\
\hline
\end{tabular}
\end{center}

For the decimal counts reported in the source, the checker verifies the
following exact integer inequalities; no logarithm or floating-point
approximation enters the checks.
\begin{center}
\begin{tabular}{r r l}
\hline
$n$&Decimal digits of $B_n$&Verified interval\\
\hline
1&1&$10^0\leq B_1<10^1$\\
10&6&$10^5\leq B_{10}<10^6$\\
100&116&$10^{115}\leq B_{100}<10^{116}$\\
1000&1928&$10^{1927}\leq B_{1000}<10^{1928}$\\
\hline
\end{tabular}
\end{center}
These intervals certify the digit counts. For compact reproducibility, the
JSON retains each value's first and last at most $32$ decimal digits and a
SHA-256 fingerprint of its full decimal expansion. The fingerprints are
identifiers; the recurrence and the integer inequalities are the mathematical
checks. No value $B_{10000}$ is computed or asserted here.

Among $0\leq n\leq13$, precisely $n=2,3,7,13$ give primes:
\[
 B_2=2,\qquad B_3=5,\qquad B_7=877,\qquad B_{13}=27644437.
\]
For these four values $v$, the integer square roots are respectively
$1,2,29,5257$. The checker verifies $r^2\leq v<(r+1)^2$ and that no integer
$2\leq d\leq r$ divides $v$. This proves primality: if $v=ab$ with
$a,b>1$, at least one factor is at most $r$. For every other index through
$13$, the output records a divisor or the fact that the value is less than
two. The checks make no primality assertion about later Bell numbers and no
claim about infinitely many Bell primes.

All these are fixed finite computation certificates. Their recurrence loops,
exact divisions, decimal comparisons, and trial-divisor searches have explicit
finite bounds and can be expanded into ground arithmetic derivations. The
Python program is not a generic proof checker, and successful execution is
not presented as a completed syntactic formalization in \ipra.

\begin{thebibliography}{99}
\addcontentsline{toc}{section}{References}

\bibitem{mathworld}
E. W. Weisstein, \emph{Bell Number}, MathWorld--A Wolfram Web Resource.
Requested statement inventory, accessed November 15, 2026.
\url{https://mathworld.wolfram.com/BellNumber.html}.

\bibitem{wikipedia}
\emph{Bell number}, Wikipedia, especially the Growth rate section.
Requested statement inventory, accessed November 15, 2026.
\url{https://en.wikipedia.org/wiki/Bell_number}.

\bibitem{spivey}
M. Z. Spivey, A generalized recurrence for Bell numbers,
\emph{Journal of Integer Sequences} \textbf{11} (2008), Article 08.2.5.
\url{https://cs.uwaterloo.ca/journals/JIS/VOL11/Spivey/spivey25.pdf}.

\bibitem{komatsu-pita}
T. Komatsu and C. Pita-Ruiz, Some formulas for Bell numbers,
\emph{Filomat} \textbf{32}(11) (2018), 3881--3889.
\url{https://doi.org/10.2298/FIL1811881K}.

\bibitem{bender-canfield}
E. A. Bender and E. R. Canfield, Log-concavity and related properties of the
cycle index polynomials, \emph{Journal of Combinatorial Theory, Series A}
\textbf{74} (1996), 57--70.
\url{https://mathweb.ucsd.edu/~ebender/reprints/93.pdf}.

\bibitem{callan-cesaro}
D. Callan, \emph{Ces\`aro's Integral Formula for the Bell Numbers (Corrected)},
October 3, 2005.
\url{https://pages.stat.wisc.edu/~callan/papersother/cesaro/cesaro.pdf}.

\bibitem{berend-tassa}
D. Berend and T. Tassa, Improved bounds on Bell numbers and on moments of
sums of random variables, \emph{Probability and Mathematical Statistics}
\textbf{30}(2) (2010), 185--205.
\url{https://www.math.uni.wroc.pl/~pms/files/30.2/Article/30.2.1.pdf}.

\bibitem{moser-wyman}
L. Moser and M. Wyman, An asymptotic formula for the Bell numbers,
\emph{Transactions of the Royal Society of Canada}, Series III,
\textbf{49}, Section 3 (1955).
\url{https://oeis.org/A000110/a000110_4.pdf}.

\bibitem{debruijn}
N. G. de Bruijn, \emph{Asymptotic Methods in Analysis}, Dover, 1981,
pp. 102--109. Historical attribution of the logarithmic expansion; the
explicit remainder used here is proved in Theorem~\ref{thm:debruijn}.

\bibitem{proveit}
V. Reshetnikov and contributors, \emph{ProveIt}, ConciseIPRA development,
commit \texttt{c0d0dc37f60aea10c0182ba4eece9eae8c7e13c6}.
\url{https://github.com/VladimirReshetnikov/ProveIt/tree/c0d0dc37f60aea10c0182ba4eece9eae8c7e13c6/Logic/PrimitiveRecursiveArithmetic/ConciseIPRA}.
The exact inspected files and Git blob identifiers are recorded in
\texttt{reference/proveit\_audit.json}.

\end{thebibliography}

\end{document}
